\documentclass[11pt]{article}

\usepackage[margin=1in]{geometry}
\usepackage{amsmath,amssymb,amsthm,mathtools}
\usepackage{array}
\usepackage{enumitem}
\usepackage[hidelinks]{hyperref}
\setlength{\emergencystretch}{3em}

\newtheorem{definition}{Definition}[section]
\newtheorem{theorem}[definition]{Theorem}
\newtheorem{corollary}[definition]{Corollary}
\newtheorem{lemma}[definition]{Lemma}
\newtheorem{proposition}[definition]{Proposition}
\newtheorem{remark}[definition]{Remark}

\DeclareMathOperator{\Ree}{Re}
\DeclareMathOperator{\Imm}{Im}
\DeclareMathOperator{\arglocal}{arg}
\DeclareMathOperator{\diag}{diag}
\DeclareMathOperator{\Tr}{Tr}
\newcommand{\C}{\mathbb{C}}
\newcommand{\R}{\mathbb{R}}
\newcommand{\N}{\mathbb{N}}
\newcommand{\Z}{\mathbb{Z}}
\newcommand{\Q}{\mathbb{Q}}
\newcommand{\T}{\mathbb{T}}
\newcommand{\conj}[1]{\overline{#1}}
\newcommand{\code}[1]{\texttt{#1}}

\title{The Double-Ended Helix and Projection Primacy:\\
A Formally Verified Three-Dimensional Carrier for Dirichlet $L$-Functions,\\
with the Generalized Riemann Hypothesis as the Completeness of its Projection}
\author{Samuel Lavery}
\date{2026-06-19}

\begin{document}
\maketitle

\begin{abstract}
We introduce a formally specified three-dimensional, double-ended \emph{Frobenius-dual helix carrier}
for Dirichlet $L$-function phasors. Uniform arclength placement on the growing-radius helix forces
$r_n\sim C\sqrt n$ (the area law), making $n^{-1/2}$ the \emph{unique} reciprocal scale for an
admissible fiber; prime multiplication acts by a similitude $w\mapsto\sqrt m\,\mathrm{wind}(m)\,w$
(\code{helixPt\_mul}, \code{frobeniusSimilitude\_mul}), and the functional equation pairs the two
chiral fibers under $s\leftrightarrow1-s$. Thus the carrier's admissible balanced-state space is
intrinsically the critical line $\Ree s=\tfrac12$ (\code{scaleBalanced\_iff},
\code{no\_native\_offaxis\_support}): off-axis abscissae are not scale-normalized and define no such
state, so the carrier has \emph{no native off-axis support} (\code{no\_native\_offaxis\_support}): its
real-height projection cannot land off the line. Whether that projection is also \emph{complete} --- whether
it exhausts the $L$-zeros --- is the \emph{projection-primacy} principle, which we prove equivalent to the
real-height exhaustion of the zeros, hence to GRH (\code{ProjectionPrimacy}, \code{projectionPrimacy\_iff\_exhausts}).
The Riemann Hypothesis is thereby reframed as the completeness of the projection: \emph{either} the more
accurate 3D model defines the 1D/2D readout in full --- and there are no off-line zeros --- \emph{or} an
off-line zero arises from a factor the 3D representation does not contain, one the carrier provably cannot
produce. This is the pivotal undecided hypothesis --- a reframing of RH as a geometric completeness
question, not a result short of it. The
fiber attached to a Dirichlet character $\chi$ is the
partial-sum sequence $\sum_{n<N}\chi(n)\,n^{-s}$; its limit is the Dirichlet $L$-function
$L(s,\chi)$ on $\Ree s>1$, and, after Abel summation, on the whole half-plane $\Ree s>0$,
the critical line included. Lifted to three dimensions the phasor splits into spinning charged
channels --- which carry the $L$-series unchanged --- and a standing neutral mass channel, conserving
each term's magnitude $n^{-1/2}$ that the planar arrow discards (\code{phasor3D\_mag3}); the on-line
accumulation converges height-by-height (\code{finiteCarrier\_tendsto\_LOnLine}), its vanishing a
cancellation of the fiber against the no-radial-drift carrier. The angular winding of the carrier is a completely multiplicative
character built from prime angles by the fundamental theorem of arithmetic, with no logarithm;
the logarithm enters only in the external bridge $\mathrm{wind}\,n\leftrightarrow n^{it}$ that
identifies the geometric winding with the analytic phasor. We record the geometric explicit
formula (the logarithmic derivative of the fiber is the winding-weighted von Mangoldt prime
field, and each simple zero is a simple pole with an explicit residue), the reality of the
completed object on the critical line (read through a projection that rescales the helix's $\pi/3$
gauge to a unit coordinate), the conjugacy of the two chiralities of the double-ended carrier, and a
determinant identity: at a conjugate crossing the two chiral eigenphases
$z=e^{-iy\log n}$ and $\conj{z}$ assemble the transverse block $\diag(z,\conj z)$, whose
determinant is $z\,\conj z=|z|^2=1$; this transverse block is in fact \emph{unitary}
(\code{frobeniusBlock\_unitary}), the orientation- and volume-preserving chiral invariant of the
transverse phase dynamics. These two chiral eigenphases are the values of a unit-modulus
eigenstate of the self-adjoint generator $D=-i\,d/dt$ (real eigenvalue, by von Neumann reality), and
under the de Branges change of variables the on-line cancellations we find are real spectral points
of a Hermite--Biehler structure function. The chiral forward/reverse pairing is a positive-definite
\emph{cup} form that completes to a Hilbert space (\code{hilbert\_completion\_exists}), on which the
prime-multiplication maps act as conformal symplectic similitudes scaling the K\"ahler form by $m$
(\code{theoremB}, with $\langle x,y\rangle=g(x,y)+i\,\Omega(x,y)$); the winding-weighted von Mangoldt
field is an operator trace equal to $-L'/L$ (\code{theoremC}), and a critical-line zero is re-expressed as a
one-dimensional helix cohomology class (\code{theoremD}). The signed channels also assemble a $2\times2$ harmonic
pencil whose determinant (equivalently Gram) rank-drop on the line is exactly an $L$-zero
(\code{harmonic\_pencil\_det\_zero\_iff\_L\_zero}), and a zero is \emph{admissible} --- represented by a
real source-height --- iff it lies on the line (\code{admissible\_iff\_re\_half}); the representability
of every zero --- completeness of the projection --- is the open question (projection primacy). That the
structure function for $\Lambda$ is Hermite--Biehler (equivalently, no off-line zeros) is the same Riemann
Hypothesis in the classical de Branges form; we use the de Branges framework only through its
unconditional on-line content and neither formalize nor rest on $\mathrm{IsHB}\,\Lambda$. The \emph{local} vanishing--eigenstate link, by contrast, is an identity
rather than a hypothesis: on the line the fiber accumulation is exactly the $\ell^2$ inner product
$\langle\psi_\gamma,\mathrm{data}\rangle$ of the sampled unit eigenstate $\psi_\gamma(t)=e^{i\gamma t}$
of $D$ with the arithmetic coefficient vector (\code{fiberCarrierFinite\_eq\_inner}), so a vanishing
point is exactly the resonance where $\psi_\gamma$ is orthogonal to the data
(\code{fiberCarrierFinite\_eq\_zero\_iff\_orthogonal}) --- exhibited at each located crossing. Only the \emph{global}
statement is hypothetical: that these eigenstates form the full spectrum of a single self-adjoint
operator with no off-line zeros. The regularized critical-line vanishings of the carrier are
\emph{exactly} the on-line zeros of $\zeta$ and of every Dirichlet $L$-function --- an identity, not a
numerical coincidence (\code{etaTwistClosed\_eq\_zero\_iff\_critical}, \code{FullFiber\_eq\_zero\_iff\_zeta}).
Every statement below is formalized in Lean~4 with
Mathlib and checked by the kernel; the axiom footprint of each named theorem is
$\{\mathrm{propext},\ \mathrm{Classical.choice},\ \mathrm{Quot.sound}\}$, with no \code{sorry}
and no additional axioms.
\end{abstract}

\tableofcontents

\section{Objects and notation}\label{sec:notation}

Let $\chi$ be a Dirichlet character modulo $q$ and $L(s,\chi)$ the associated $L$-function;
$\zeta$ is the Riemann zeta function and $\Lambda(s)=\code{completedRiemannZeta}(s)$ its completed
form, normalized so that $\Lambda(1-s)=\Lambda(s)$. The \emph{critical line} is the set
$s=\tfrac12+iy$, $y\in\R$. We write $\T=\{z\in\C:|z|=1\}$ for the unit circle (Mathlib
\code{Circle}) and $\conj{\,\cdot\,}$ for complex conjugation (\code{starRingEnd}~$\C$). For an
arithmetic weight $\chi$ the \emph{phasor} attached to the integer $n\ge 1$ on the vertical line
$s=\sigma+iy$ is $\chi(n)\,n^{-s}$.

\medskip
\noindent\emph{The abscissa $\tfrac12$.} Every $\tfrac12$ below is one of two \emph{derived}
constants, neither of which assumes where the zeros lie. As the carrier's \emph{scale-critical
exponent} it is the unique $\sigma$ at which the phasor amplitude $n^{-\sigma}$ balances the emergent
area-law radius $r_n\sim\sqrt n$ (Theorem~\ref{thm:scalecrit}); this fixes the amplitude $n^{-1/2}$
and the abscissa of the line the fiber rides, from the carrier geometry alone. As the \emph{reflection
axis} $\Ree s=\tfrac12$ it is the fixed line of $s\mapsto 1-s$, on which $\conj s=1-s$; with the
functional equation $\Lambda(1-s)=\Lambda(s)$ this is what makes the completed object real on the line
(Theorem~\ref{thm:real}) and pairs the two chiralities (Theorem~\ref{thm:bothsides}). The Riemann
Hypothesis --- that the \emph{zeros} $\rho$ satisfy $\Ree\rho=\tfrac12$ --- is a third, distinct
statement, never assumed here: the on-line results below are either unconditional in $y$ or
conditional on a zero being present at the stated height, and they assert nothing about off-line
zeros.

The classical analytic theory of $\zeta$ and Dirichlet $L$-functions --- the Euler product, the
functional equation, and the von Mangoldt explicit formula
\cite{Riemann1859,Edwards1974,Titchmarsh1986,Davenport2000,IwaniecKowalski2004,MontgomeryVaughan2007}
--- enters here through the corresponding Mathlib lemmas.

All named results carry the identifier of the corresponding Lean declaration, e.g.\
\code{frobenius\_conjugate\_det\_one}; Section~\ref{sec:formal} records the verification.

\section{The $\pi/3$ helix carrier}\label{sec:carrier}

\subsection{The space curve and its arclength}

\begin{definition}[Carrier curve; \code{Geometry.helix}, \code{helixVel}, \code{speed}]\label{def:helix}
For pitch $p$ and radial rate $r$, the carrier is the growing-radius helix
\[
  \gamma(k) \;=\; \bigl(r\,k\cos 2\pi k,\; r\,k\sin 2\pi k,\; p\,k\bigr),\qquad k\in\R .
\]
Its velocity is $\gamma'(k)=\bigl(r\cos2\pi k-2\pi r k\sin2\pi k,\ r\sin2\pi k+2\pi r k\cos2\pi k,\ p\bigr)$,
and the squared speed is
\[
  \|\gamma'(k)\|^2 \;=\; p^2+r^2+(2\pi r k)^2 .
\]
The \emph{climber} is $k(y)=e^{y}/p$, so the height is $z=p\,k(y)=e^{y}$ and the cylindrical radius
is $R(y)=r\,k(y)=(r/p)e^{y}$.
\end{definition}

\begin{remark}[Constant height law: $\Delta z=1$ per integer $n$, not per phasor; \code{Geometry.heightLaw}, \code{heightLaw\_step}, \code{resonance\_height}, \code{activeFraction\_eq\_totient}, \code{carrierRadius\_growth\_div\_sqrt\_tendsto}]\label{rem:height}
Discretely, the carrier assigns integer $n$ the height $z_n=n$: a constant $\Delta z=1$
\emph{per integer $n$} (\code{heightLaw\_step}: $\mathrm{heightLaw}(n+1)-\mathrm{heightLaw}(n)=1$)---every
slot, whether or not its arithmetic weight $\chi(n)$ is nonzero.
A neutral slot ($\chi(n)=0$) carries no arrow yet still occupies its place and still advances the
height. Counting instead per \emph{active} phasor would climb by the active fraction
$\varphi(q)/q$ per integer (\code{activeFraction\_eq\_totient}: only the $\varphi(q)$ unit residues carry a
spinning arrow) and land the resonance at $(\varphi(q)/q)\,e^{\gamma}$, short by exactly
that fraction. Since phasor $n$ resonates at $y=\log n$, integer $n$ sits at $z_n=n=e^{\log n}$---the
height of its own resonant frequency (\code{resonance\_height})---so a zero at frequency $\gamma$ (resonant term
$n^\ast=e^{\gamma}$) lands at height $z=e^{\gamma}$, for every $L$-function (it rides only on
$\mathrm{spin}=\log n$; the located on-line vanishings are exactly the carrier's resonant heights,
\code{online\_zero\_eq\_carrierZeros}, each at a unique height, \code{online\_zero\_unique\_height}). This is the discrete form of the climber height $z=p\,k(y)=e^{y}$ above.
The radial gap per turn is $L$-dependent, $g=e^{\,(q\bmod{})}$ (eta $e^2$, $\chi_3$ $e^3$, \dots), so the
area-law radius specializes to $r_n/\sqrt n\to\sqrt{e^{q}/3}$ (\code{carrierRadius\_growth\_div\_sqrt\_tendsto}):
the same $\sqrt n$ shape for every $L$, only the constant differs --- consistent with the
gauge-independence of the scale-critical exponent (Theorem~\ref{thm:scalecrit}).
\end{remark}

\begin{definition}[Arclength; \code{Geometry.arclength}, \code{arclengthClosed}]\label{def:arclength}
The arclength of $\gamma$ from $0$ to $k$ is
$\displaystyle S(k;p,r)=\int_0^k\sqrt{p^2+r^2+(2\pi r t)^2}\,dt$.
\end{definition}

\begin{theorem}[Closed-form arclength; \code{arclength\_closed\_form}, \code{arclength\_r\_zero}]\label{thm:arclength}
For $r>0$,
\[
  S(k;p,r) \;=\; \frac{k}{2}\sqrt{p^2+r^2+4\pi^2 r^2 k^2}
   \;+\;\frac{p^2+r^2}{4\pi r}\,\operatorname{arsinh}\!\Bigl(\frac{2\pi r k}{\sqrt{p^2+r^2}}\Bigr),
\]
and $S(k;p,0)=p\,k$ for $p\ge 0$.
\end{theorem}

\begin{proof}[Proof sketch]
Differentiate the stated antiderivative and check it equals the integrand
$\sqrt{p^2+r^2+4\pi^2r^2k^2}$ pointwise, then apply the fundamental theorem of calculus on
$[0,k]$; the constant-pitch case is $\int_0^k p\,dt$.
\end{proof}

\subsection{Uniform integer placement and the area law}

\begin{definition}[Spacing, index, integer angle; \code{Geometry.Delta}, \code{Nindex}, \code{spinAngle}]\label{def:placement}
The fixed spacing is $\Delta=\pi/3$. The continuous geometric index is $N(y)=S(k(y);p,r)/\Delta$,
so $N(y)=(3/\pi)\,S(k(y);p,r)$ (\code{Nindex\_eq}). The integer angular coordinate is
$s_n=n\cdot(\pi/3)$.
\end{definition}

\begin{theorem}[Uniform $\pi/3$ spacing; \code{carrier\_spacing}]\label{thm:spacing}
For every $n\in\N$, $\;s_{n+1}-s_n=\pi/3$.
\end{theorem}

\begin{theorem}[Area-law bound; \code{arclengthClosed\_lower\_bound}, \code{arclengthClosed\_upper\_bound}]\label{thm:areabound}
For $r>0$ and $k\ge 0$,
\[
  \pi r\,k^2 \;\le\; S(k;p,r) \;\le\; \pi r\,k^2 + 2k\sqrt{p^2+r^2}.
\]
\end{theorem}

\begin{definition}[Wound integer placement; \code{Geometry.windParameter}, \code{windIntegerSite}, \code{exists\_unique\_windParameter}]\label{def:wound}
Integer $n$ sits at arclength $s_n=n\Delta$ on the unwound line (Definition~\ref{def:placement}).
\emph{Winding the line onto the helix} places it at the unique parameter $k_n\ge 0$ with
$S(k_n;p,r)=n\Delta$ (\code{exists\_unique\_windParameter}; the map $s\mapsto k_n$ is
\code{windParameter}). The wound site is $\mathrm{windIntegerSite}(n)=\gamma(k_n)$, with cylindrical
radius $r_n=r\,k_n$ (\code{windIntegerSite\_cyl\_radius}).
\end{definition}

\begin{theorem}[Emergent radius / area law; \code{emergent\_radius\_sq\_tendsto}, \code{windIntegerSite\_radius\_sq\_tendsto}]\label{thm:arealaw}
For the wound placement $S(k_n;p,r)=n\Delta$ of Definition~\ref{def:wound},
\[
  \frac{r_n^{\,2}}{n}\;=\;\frac{(r\,k_n)^2}{n}\;\xrightarrow[\;n\to\infty\;]{}\;\frac{r\Delta}{\pi},
  \qquad\text{equivalently}\qquad r_n\;\sim\;\sqrt{\tfrac{r\Delta}{\pi}}\;\sqrt n,
\]
so the enclosed area $\pi r_n^{\,2}\sim r\Delta\,n$ is linear in $n$. The factor $\sqrt n$ is thus
\emph{derived from the arclength geometry}: the bound of Theorem~\ref{thm:areabound} forces
$S(k)\propto k^2$, hence $k_n\propto\sqrt n$ and $r_n=r\,k_n\propto\sqrt n$ --- it is not posited.
\end{theorem}

\begin{remark}[Unit gauge; \code{windIntegerSite\_radius\_sq\_tendsto\_unit\_gauge}]\label{rem:gauge}
The area-law constant is $r\Delta/\pi$; with the native spacing $\Delta=\pi/3$ it equals $r/3$, so
$r_n\sim\sqrt n$ \emph{exactly} in the gauge $r\Delta=\pi$, i.e.\ $r=3$
(\code{windIntegerSite\_radius\_sq\_tendsto\_unit\_gauge}); in the alternative gauge $r=1$ the
constant is $1/\sqrt3$. In this unit gauge the carrier point
$\mathrm{helixPt}(n)=\sqrt n\cdot\mathrm{wind}(n)$ (Definition~\ref{def:wind}) has modulus
$\|\mathrm{helixPt}(n)\|=\sqrt n$ (\code{HelixLogFree.norm\_helixPt}), matching the emergent radius
$r_n$ asymptotically (\code{helixPt\_radius\_matches\_areaLaw}: the ratio
$\|\mathrm{helixPt}(n)\|/r_n\to 1$ in this gauge); the winding contributes unit modulus
(Theorem~\ref{thm:windmul}).
\end{remark}

\subsection{The scale-critical exponent}\label{sub:scalecrit}

The area law gives more than the radius $\sqrt n$: it pins the critical exponent. Write
$r_n=\mathrm{carrierRadius}(n)$ for the wound radius of Definition~\ref{def:wound}
(\code{Geometry.carrierRadius}), and pair it with a candidate weight $n^{-\sigma}$, $\sigma\in\R$.

\begin{theorem}[The critical exponent is scale-critical; \code{sigma\_half\_is\_scale\_critical}]\label{thm:scalecrit}
For $\sigma\in\R$, the scale-balance product $n^{-\sigma}\,r_n$ has a positive limit if and only if
$\sigma=\tfrac12$:
\[
  \bigl(\exists L>0:\ n^{-\sigma}\,r_n\xrightarrow[\;n\to\infty\;]{}L\bigr)
  \qquad\Longleftrightarrow\qquad \sigma=\tfrac12 .
\]
Concretely, $n^{-\sigma}r_n\to 0$ for $\sigma>\tfrac12$ (the weight decays faster than the radius
grows), $n^{-\sigma}r_n\to+\infty$ for $\sigma<\tfrac12$ (the weight outruns the radius), and
$n^{-\sigma}r_n\to\sqrt{r\Delta/\pi}>0$ for $\sigma=\tfrac12$. Thus $\tfrac12$ is the unique exponent at
which the weight reciprocates the emergent radius. It is \emph{gauge-independent}: the gauge $r\Delta$
sets only the value of the limit, never which $\sigma$ is critical.
\end{theorem}

\begin{proof}[Proof sketch]
By Theorem~\ref{thm:arealaw}, $r_n/\sqrt n\to\sqrt{r\Delta/\pi}=:c>0$. Factor
$n^{-\sigma}r_n=(r_n/\sqrt n)\cdot n^{1/2-\sigma}$; the first factor tends to $c$, and the second tends
to $0$, to $1$, or to $+\infty$ according as $\sigma>\tfrac12$, $\sigma=\tfrac12$, or $\sigma<\tfrac12$.
Only $\sigma=\tfrac12$ yields a positive finite limit; uniqueness of limits closes both directions.
\end{proof}

\begin{remark}\label{rem:scalecrit}
This is the geometric origin of the critical line: the $\tfrac12$ of $\Ree s=\tfrac12$ is the $\tfrac12$
of $r_n\sim\sqrt n=n^{1/2}$. It locates the line on which the phasor weights are scale-normalized to
$O(1)$, so that the accumulation is governed by the winding (the phase cancellation) rather than the
amplitude --- which is exactly the mechanism by which the vanishing points arise
(Section~\ref{sec:strip}). It is a statement about the carrier alone; that the zeros of $\zeta$ and
$L$ then \emph{lie} on this line is the separate analytic content of the functional equation
(Theorem~\ref{thm:bothsides}), not of the area law.
\end{remark}

\subsection{The mod-$6$ bucket structure}

Because $6\cdot(\pi/3)=2\pi$, the integer angle $s_n$ is $6$-periodic on the circle.

\begin{theorem}[Bucket periodicity and counts; \code{spin\_phasor\_period6}, \code{bucket\_count}]\label{thm:buckets}
$\exp(i\,s_{n+6})=\exp(i\,s_n)$ for all $n$, and for $a<6$ exactly $N$ of the integers in
$[0,6N)$ satisfy $n\equiv a\pmod 6$. Thus the integers collapse onto six angular directions,
each receiving an equal share.
\end{theorem}

\begin{remark}[Sixth roots of unity; Eisenstein integers]\label{rem:eisenstein}
The integer-angle phasor $\exp(i\,s_n)=\exp(i\,n\pi/3)$ takes values in the sixth roots of unity
$\mu_6=\{e^{ik\pi/3}:0\le k<6\}$, cycling through them as $n$ runs mod $6$; the primitive sixth root
$\zeta=e^{i\pi/3}$ ($\zeta^2=\zeta-1$, the cyclotomic relation $\Phi_6$) generates $\mu_6$. The ring
$\Z[\zeta]$ is the ring of Eisenstein integers \cite{IrelandRosen1990} --- a hexagonal (triangular)
lattice whose unit group
is exactly $\mu_6$. So the $\pi/3$ gauge is the Eisenstein symmetry, and the six buckets of
Theorem~\ref{thm:buckets} are its units.
\end{remark}

\begin{theorem}[Exact finite monodromy of the $\pi/3$ gauge; \code{pi3\_finite\_monodromy}, \code{unit\_not\_finite\_monodromy}, \code{circleMonodromy\_iff\_rat}]\label{thm:monodromy}
Read the carrier phase after $k$ cells modulo a turn, $C_H(k)=k\,H \pmod{2\pi}$; a cell unit $H$ has
\emph{exact finite circle monodromy} when the carrier returns to its origin after some positive number
of cells, i.e.\ $k\,H=2\pi m$ for integers $k\ge1$, $m\ge1$. The native gauge $H=\pi/3$ closes
\emph{exactly} after six cells, $6\cdot(\pi/3)=2\pi$ ($k=6$, $m=1$; \code{pi3\_finite\_monodromy}),
whereas the unit gauge $H=1$ never closes (\code{unit\_not\_finite\_monodromy}: it would force $\pi\in\Q$).
In general $H$ has exact finite monodromy \emph{iff} $H/(2\pi)\in\Q_{>0}$
(\code{circleMonodromy\_iff\_rat}). Thus $\pi/3$ is distinguished not only as the Eisenstein unit angle
(Remark~\ref{rem:eisenstein}) but as a \emph{finite-monodromy} gauge --- the exact closing condition
underlying the six buckets of Theorem~\ref{thm:buckets}.
\end{theorem}

\begin{remark}[A closed carrier for an open winding]\label{rem:closedvsopen}
The exactness of this closure is a genuine structural distinction from the object the carrier serves.
The analytic phasor winding $n^{-it}=e^{-it\log n}$ rides incommensurate frequencies --- the $\log p$ are
$\Q$-linearly independent --- so it never returns: the standard presentation of $L$ has no closed angular
structure. The $\pi/3$ carrier, by contrast, closes \emph{exactly} after six cells, and only the
rational-angle gauges close at all (Theorem~\ref{thm:monodromy}). The carrier thus supplies a closed,
finite, commensurate angular home --- the Eisenstein/$\mu_6$ hexagonal lattice
(Remark~\ref{rem:eisenstein}) --- for an arithmetic that has none in its analytic form. (The cell
enters the readout as a unit-modulus factor, divided out to recover the $L$-value in Section~\ref{sec:strip}.)
\end{remark}

\section{The log-free winding and the bridge}\label{sec:winding}

\subsection{The fundamental-theorem-of-arithmetic winding}

\begin{definition}[Winding angle and winding; \code{HelixLogFree.windAngle}, \code{wind}, \code{helixPt}]\label{def:wind}
Fix a prime-angle assignment $\theta$ (in Lean, a function $\theta:\N\to\R$; only its values on primes
enter).
The \emph{winding angle} is the completely additive extension read off the prime factorization,
\[
  \Theta(n) \;=\; \sum_{p^{e}\,\|\,n} e\,\theta(p)
  \;=\; \code{n.factorization.sum}\ (\lambda\, p\, e.\; e\cdot\theta(p)),
\]
and the \emph{winding} is $\mathrm{wind}(n)=\exp(i\,\Theta(n))\in\T$. The carrier point is
$\mathrm{helixPt}(n)=\sqrt n\cdot \mathrm{wind}(n)$. No logarithm occurs in $\Theta$.
\end{definition}

\begin{theorem}[Multiplicative character; \code{windAngle\_mul}, \code{wind\_mul}, \code{wind\_one}]\label{thm:windmul}
For $m,n\neq 0$, $\;\Theta(mn)=\Theta(m)+\Theta(n)$ and $\mathrm{wind}(mn)=\mathrm{wind}(m)\,\mathrm{wind}(n)$,
with $\mathrm{wind}(1)=1$. Thus $\mathrm{wind}:\N\to\T$ is a completely multiplicative character
\cite{Apostol1976}.
\end{theorem}

\begin{proof}[Proof sketch]
$\Theta$ is a sum over $\code{Nat.factorization}$; additivity is $\code{Nat.factorization\_mul}$
(the fundamental theorem of arithmetic). Exponentiating turns additivity into multiplicativity on
$\T$ via $\code{Circle.exp\_add}$.
\end{proof}

\subsection{The bridge $\mathrm{wind}\,n\leftrightarrow n^{it}$}

The single place a logarithm enters is the external identification of the geometric winding with the
analytic phasor.

\begin{theorem}[Logarithm is FTA-additive; \code{real\_log\_eq\_factorization\_sum}]\label{thm:logfta}
For $n\neq 0$, $\;\log n=\sum_{p^{e}\,\|\,n} e\,\log p$.
\end{theorem}

\begin{theorem}[The bridge; \code{windAngle\_glog}, \code{wind\_glog\_eq\_cpow}]\label{thm:bridge}
With the prime-angle assignment $\theta(p)=\gamma\log p$, one has
$\Theta(n)=\gamma\log n$ and hence
\[
  \mathrm{wind}(n) \;=\; n^{\,i\gamma}\qquad(n\ge 1).
\]
\end{theorem}

\begin{proof}[Proof sketch]
Substitute Theorem~\ref{thm:logfta} into $\Theta$ to get $\Theta(n)=\gamma\log n$; then
$\exp(i\gamma\log n)=n^{i\gamma}$ by the principal-branch definition of $n^{i\gamma}$ for the
positive real base $n$.
\end{proof}

\section{The phasor fiber and its accumulation}\label{sec:fiber}

\subsection{Phasor decomposition}

\begin{theorem}[Vertical-line phasor and magnitude; \code{cpow\_vertical\_line\_phasor}, \code{norm\_cpow\_vertical\_line}]\label{thm:phasor}
For a positive real base $x$ and $s=\sigma+iy$,
\[
  x^{-s} \;=\; x^{-\sigma}\,\exp\!\bigl(-(y\log x)\,i\bigr),\qquad
  \bigl|x^{-s}\bigr| \;=\; x^{-\sigma}.
\]
In particular, on the critical line $n^{-(1/2+iy)}=n^{-1/2}\,\exp(-(y\log n)i)$ with magnitude
$n^{-1/2}$ for every $y$ (\code{cpow\_critical\_line}, \code{norm\_cpow\_critical\_line}).
\end{theorem}

\begin{definition}[The spin; \code{LFunctionPhasor.spin}]\label{def:spin}
The \emph{spin} of the integer $n$ at ordinate $y$ is the unit-modulus factor
$\mathrm{spin}(y,n)=\exp(-(y\log n)\,i)$.
\end{definition}

\begin{theorem}[The fiber rides the carrier; \code{fiber\_rides\_carrier}]\label{thm:rides}
Write $\mathrm{wind}_{-t\log}$ for the winding of Definition~\ref{def:wind} with the assignment
$\theta(p)=-t\log p$, so $\mathrm{wind}_{-t\log}(n)=n^{-it}$ by the bridge (Theorem~\ref{thm:bridge}).
For $n\ge 1$, the carrier magnitude $n^{-1/2}$ --- the reciprocal of the emergent radius $\sqrt n$ ---
times this winding is the Dirichlet phasor:
\[
  n^{-1/2}\cdot\mathrm{wind}_{-t\log}(n) \;=\; n^{-(1/2+it)} .
\]
\end{theorem}

\begin{remark}[The fiber's exponent is scale-critical; \code{critical\_exponent\_is\_scale\_critical}]\label{rem:ridescrit}
The exponent $\tfrac12$ here is not a chosen constant. By Theorem~\ref{thm:scalecrit} --- re-exported at
the unit-gauge carrier $r=3$ as \code{critical\_exponent\_is\_scale\_critical} --- it is the unique
$\sigma$ for which the fiber amplitude $n^{-\sigma}$ scale-balances the area-law radius $r_n$. So the
critical line $\Ree s=\tfrac12$ that the fiber rides is the scale-balance locus of the carrier, forced
by the geometry rather than inserted.
\end{remark}

\subsection{Accumulation to $L$}

\begin{definition}[Finite phasor carrier; \code{DirichletPhasorCarrier.finiteCarrier}, \code{phasorTerm}]\label{def:finite}
$G_{\chi,N}(s)=\sum_{n<N}\chi(n)\,n^{-s}$, the partial sums of the carrier-riding phasors.
\end{definition}

\begin{theorem}[Accumulation on $\Ree s>1$; \code{fiber\_accumulates\_to\_L}, \code{finiteCarrier\_tendsto\_LFunction}]\label{thm:accum}
For every Dirichlet character $\chi$ and $\Ree s>1$, $\;G_{\chi,N}(s)\to L(s,\chi)$ as $N\to\infty$.
\end{theorem}

\section{The three-channel three-dimensional phasor}\label{sec:phasor3d}

The planar phasor $\chi(n)\,n^{-\sigma}\exp(-(y\log n)i)$ of Section~\ref{sec:fiber} is a 2D arrow
whose length is $\|\chi(n)\|\,n^{-\sigma}$; on the \emph{neutral} channel $\chi(n)=0$ it collapses to
$0$ and the magnitude $n^{-\sigma}$ is lost. The carrier, by contrast, places \emph{every} slot $n$
(Remark~\ref{rem:height}), neutral or not. The reconciliation lifts each term to a genuine vector in
$\R^3=(\text{spin plane }\R^2)\times(\text{mass axis }\R)$, modeled as $\C\times\R$.

\begin{definition}[3D phasor and its three channels; \code{Phasor3D.phasor3D}, \code{channel}]\label{def:phasor3d}
\[
  \mathrm{phasor3D}(\chi,\sigma,y,n)=\Bigl(\,\underbrace{\chi(n)\,n^{-\sigma}\exp(-(y\log n)i)}_{\text{spin plane}},\
  \underbrace{(1-\|\chi(n)\|)\,n^{-\sigma}}_{\text{mass axis}}\,\Bigr).
\]
The charged channels $\chi(n)=\pm1$ place the rotating arrow $\pm n^{-\sigma}\exp(-(y\log n)i)$ in the
spin plane with zero mass; the neutral channel $\chi(n)=0$ places the real amplitude $n^{-\sigma}$ on
the mass axis with no spin.
\end{definition}

\begin{theorem}[The spin plane carries $L$ unchanged; \code{phasor3D\_plane\_tsum}]\label{thm:plane}
The spin-plane (complex) coordinate of $\mathrm{phasor3D}$ is exactly the planar phasor of
Section~\ref{sec:fiber}, so its accumulation is the same Dirichlet $L$-series. The 3D lift changes
nothing in the analytic representation --- it only records what the plane discards.
\end{theorem}

\begin{theorem}[Magnitude conservation across all channels; \code{phasor3D\_mag3}]\label{thm:mag3}
For a quadratic character the full $\R^3$ Euclidean length of every term is exactly $n^{-\sigma}$, on
\emph{all three} channels including the neutral one --- where the planar length $\|\chi(n)\|\,n^{-\sigma}$
vanished. The magnitude lost from the plane is precisely the mass stored on the axis.
\end{theorem}

\begin{theorem}[Standing mass vs.\ travelling spin; \code{phasor3D\_neutral\_no\_spin}, \code{phasor3D\_charged\_spin}, \code{resonantHeight\_eq\_exp\_resonantFreq}]\label{thm:standing}
The neutral term is independent of the frequency $y$ (\code{phasor3D\_neutral\_no\_spin}): a
\emph{standing} (resonant) mode, not a travelling one, storing positive amplitude $n^{-\sigma}$ on the
mass axis (\code{phasor3D\_neutral\_mass\_pos}) while contributing $0$ to the spin-plane $L$-series
(\code{phasor3D\_neutral\_plane\_zero}). The charged terms rotate at rate $\log n$ as $y$ advances
(\code{phasor3D\_charged\_spin}). The standing mass resonates at frequency
$\mathrm{resonantFreq}(n)=\log n$ and height $\mathrm{resonantHeight}(n)=n=e^{\log n}$
(\code{resonantHeight\_eq\_exp\_resonantFreq}), matching the carrier height law $z_n=n=e^{y_n}$ of
Remark~\ref{rem:height}.
\end{theorem}

\begin{remark}[Three channels, one decomposition; \code{phasor3D\_three\_channel\_form}]\label{rem:threechannel}
Every term lies in exactly one channel --- positive ($\chi(n)=+1$), negative ($\chi(n)=-1$), neutral
($\chi(n)=0$) --- and is determined by it (\code{phasor3D\_three\_channel\_form}). The signed pair
(positive/negative) is the spinning content the $L$-series sees; the neutral channel is the mass the
planar model omits. This positive/negative/neutral split is the channel structure built into the
harmonic pencil of Section~\ref{sec:pencil}.
\end{remark}

\section{Strip extension: accumulation onto the critical line}\label{sec:strip}

The accumulation converges beyond the region of absolute convergence, down to $\Ree s>0$.

\begin{theorem}[Dirichlet strip extension; \code{dirichlet\_strip\_tendsto\_LFunction}]\label{thm:dstrip}
For a Dirichlet character $\chi\neq 1$ modulo $q$ and $\Ree s>0$,
\[
  \sum_{n<N}\chi(n)\,n^{-s}\;\xrightarrow[N\to\infty]{}\;L(s,\chi).
\]
\end{theorem}

\begin{theorem}[Eta strip extension; \code{eta\_strip\_tendsto}]\label{thm:etastrip}
For $\Ree s>0$ and $s\neq 1$,
\[
  \sum_{n<N}(-1)^{\,n+1}\,n^{-s}\;\xrightarrow[N\to\infty]{}\;\bigl(1-2^{\,1-s}\bigr)\zeta(s).
\]
\end{theorem}

\begin{proof}[Proof sketch]
Bucket cancellation gives a uniform bound on the character sums $\sum_{n<N}\chi(n)$; Abel summation
by parts \cite{Apostol1976,Titchmarsh1986} then yields convergence of $\sum\chi(n)n^{-s}$ on
$\Ree s>0$, and the limit agrees with
$L(s,\chi)$ on $\Ree s>1$ and hence everywhere on the connected strip by the identity theorem. For
the alternating weight the factor $1-2^{1-s}$ is the even-term correction relating $\eta$ and $\zeta$.
\end{proof}

\begin{theorem}[Continuous reach to the line; \code{continuous\_model\_zeta}]\label{thm:cont}
For every $y\in\R$, at $s=\tfrac12+iy$,
\[
  \sum_{n<N}(-1)^{\,n+1}\,n^{-s}\;\longrightarrow\;\bigl(1-2^{\,1-s}\bigr)\zeta\!\bigl(\tfrac12+iy\bigr),
\]
and the limit is $0$ if and only if $\zeta(\tfrac12+iy)=0$.
\end{theorem}

\begin{proof}[Proof sketch]
Theorem~\ref{thm:etastrip} at $s=\tfrac12+iy$ (using $\Ree s=\tfrac12>0$ and $s\neq1$). On the line
the factor $1-2^{1-s}$ is nonzero (\code{correction\_factor\_ne\_zero\_on\_critical\_line}), so the
product vanishes exactly when $\zeta(\tfrac12+iy)$ does (\code{etaTrivial\_eq\_zero\_iff}).
\end{proof}

\subsection{Per-height convergence: the cancellation is in the fiber, not the carrier}\label{sub:perheight}

The accumulation converges height-by-height, and its vanishing is a property of the \emph{fiber},
never of the carrier. The two objects are distinct. The carrier
$\mathrm{helixPt}(n)=\sqrt n\,\mathrm{wind}(n)$ has modulus exactly $\sqrt n$, independent of the
ordinate $y$ and the character $\chi$ (\code{carrier\_norm\_height\_independent}), and never vanishes
off the origin (\code{carrier\_ne\_zero}): it is a no-radial-drift source. The fiber paired against it
is the unit-modulus spin (\code{fiber\_spin\_unimodular}) times the arithmetic datum
$\chi(n)\,n^{-1/2}$; attaching it only \emph{rotates} the carrier and cannot alter the $\sqrt n$
radial profile. Hence every vanishing of an accumulation is a fiber/pairing cancellation
(orthogonality, Proposition~\ref{prop:ortho}), never a collapse of the carrier.

\begin{theorem}[Per-height channels; \code{finiteCarrier\_critical\_tendsto}, \code{eta\_critical\_tendsto}, \code{perHeight\_channel\_nonprincipal}, \code{perHeight\_channel\_zeta}]\label{thm:perheight}
For every ordinate $y\in\R$:
\begin{enumerate}[label=\textup{(\roman*)}]
\item for non-principal $\chi$ the raw weighted partial sums $\sum_{n<N}\chi(n)\,n^{-(1/2+iy)}$
converge to $L(\tfrac12+iy,\chi)$ (\code{finiteCarrier\_critical\_tendsto}), so the convergent channel's
nonvanishing factor is $1$;
\item for the principal character (and $\zeta$) the raw series diverges, but the alternating (eta)
channel $\sum_{n<N}(-1)^{n+1}n^{-s}$ converges to $(1-2^{1-s})\zeta(s)$, whose correction factor is
nonzero on the line (\code{eta\_critical\_tendsto}, \code{zeta\_eta\_factor\_ne\_zero}).
\end{enumerate}
In both cases the channel value vanishes exactly when the $L$-value does
(\code{LFunction\_critical\_eq\_zero\_iff\_pairing\_tendsto\_zero},
\code{riemannZeta\_critical\_eq\_zero\_iff\_eta\_tendsto\_zero}), the on-line limit
$y\mapsto L(\tfrac12+iy,\chi)$ is continuous (\code{continuous\_LOnLine}), and it is the analytic limit
of the finite carrier (\code{finiteCarrier\_tendsto\_LOnLine}).
\end{theorem}

This is the precise, defensible form of ``every zero on the line is subject to the cancellation
machine'': it asserts nothing about \emph{where} the zeros are, only that each on-line zero is exactly
a fiber cancellation of the convergent channel.

\subsection{A fully convergent worked fibre for $\zeta$}\label{sub:fullfiber}

For the Riemann zeta channel the convergence above is realized as an explicit, quantitatively
controlled Cauchy sequence. Place the fibre at the critical-line point $\mathrm{sFiber}(T)$, which has
real part $\tfrac12$ by construction (\code{sFiber\_re}).

\begin{theorem}[The $\eta$-fibre converges on the line; \code{Pi3Fiber.PartialFiber\_tendsto\_FullFiber}, \code{PartialFiber\_cauchySeq}, \code{FullFiber\_eq\_zero\_iff\_zeta}]\label{thm:fullfiber}
The alternating partial sums
$\mathrm{PartialFiber}(N,T)=\sum_{n<N}(-1)^{n}\,(n+1)^{-\mathrm{sFiber}(T)}$ form a Cauchy sequence
(\code{PartialFiber\_cauchySeq}) converging to $\mathrm{FullFiber}(T)=\eta(\mathrm{sFiber}(T))$
(\code{PartialFiber\_tendsto\_FullFiber}, with vanishing tail \code{TailFiber\_tendsto\_zero}); and it
vanishes \emph{iff} $\zeta(\mathrm{sFiber}(T))=0$ (\code{FullFiber\_eq\_zero\_iff\_zeta}). The phasor form
of the partial sum exhibits the spin $\exp(-(HT\log(n+1))i)$ on the amplitude $(n+1)^{-1/2}$
(\code{PartialFiber\_eq\_phasor}).
\end{theorem}

\begin{proof}[Proof sketch]
The summand differences are controlled by the consecutive-power bound
$\bigl\|a^{-w}-(a+1)^{-w}\bigr\|\le\|w\|\,a^{-\Ree w-1}$ for $a\ge1$, $\Ree w>0$
(\code{cpow\_neg\_consecutive\_norm\_le}), giving summability of the bracketed series
(\code{pairedTerm\_summable}) and the Cauchy property on $\Ree s>0$. The bracketed sum agrees with
$\eta$ for $\Ree s>1$ and extends by the identity theorem across the preconnected region
$\{\Ree s>0\}\setminus\{1\}$ (\code{pairedEta\_eq\_etaTrivial},
\code{rightHalfPlane\_sdiff\_one\_isPreconnected}, \code{sFiber\_ne\_one}).
\end{proof}

\begin{remark}[Exact accumulation, not approximate evaluation]\label{rem:exactvsapprox}
The critical-line behavior of $L$ is classically reached by \emph{approximate} means --- the approximate
functional equation and the Riemann--Siegel formula --- which represent $L(\tfrac12+it)$ by truncated
sums with a controlled error term. The accumulation here is \emph{exact}: the regularized carrier
converges to $L$ on the line (Theorems~\ref{thm:perheight},~\ref{thm:fullfiber}) and vanishes \emph{iff}
$L$ does (Theorem~\ref{thm:located}), an identity with no error term, machine-checked. Correspondingly the
vanishing is an exact \emph{event} --- the sampled eigenwave is exactly $\ell^2$-orthogonal to the
arithmetic data (Proposition~\ref{prop:ortho}) --- rather than an approximate minimum of $|L|$. The
exactness is what lets the carrier's geometric and spectral structure transfer to $L$ rigorously rather
than heuristically; it is the foundation the rest of the development stands on. (The underlying
convergence is the classical Abel/$\eta$ summation on $\Ree s>0$; what is new is the exact, machine-checked
geometric carrier built on it.)
\end{remark}

\section{The geometric explicit formula}\label{sec:explicit}

The logarithmic derivative of the fiber is the prime-power accumulation weighted by the winding.

\begin{theorem}[Prime field; \code{explicit\_formula\_prime\_field}, \code{dirichlet\_logDeriv\_eq\_tsum}]\label{thm:primefield}
For $\Ree s>1$,
\[
  -\frac{L'(s,\chi)}{L(s,\chi)} \;=\; \sum_{n}\chi(n)\,\Lambda(n)\,n^{-s},
\]
where $\Lambda$ is the von Mangoldt function. The zeros of $L(\cdot,\chi)$ are the poles of the
meromorphic continuation of this object.
\end{theorem}

\begin{definition}[Residue harmonic; \code{Residue.residueHarmonic}]\label{def:resharm}
For $x>0$ and $\gamma\in\R$, write $\rho=\tfrac12+i\gamma$ and set
$\;\mathrm{residueHarmonic}(x,\gamma)=-x^{\rho}/\rho$.
\end{definition}

\begin{theorem}[Zeros located as poles; \code{zeta\_zero\_located\_as\_pole}, \code{L\_zero\_located\_as\_pole}]\label{thm:pole}
Let $\rho=\tfrac12+i\gamma$ be a zero of $L(\cdot,\chi)$ at which $L'(\rho,\chi)\neq 0$ (a simple
zero). Then $\rho$ is a simple pole of the explicit-formula kernel $-(L'/L)(s)\,x^{s}/s$, with
\[
  \lim_{s\to\rho}(s-\rho)\,\Bigl(-\tfrac{L'(s,\chi)}{L(s,\chi)}\cdot\tfrac{x^{s}}{s}\Bigr)
   \;=\; \mathrm{residueHarmonic}(x,\gamma)\;=\;-\frac{x^{\rho}}{\rho}.
\]
The same statement holds verbatim for $\zeta$.
\end{theorem}

\begin{proof}[Proof sketch]
On the punctured neighborhood of $\rho$ in $\{s\neq 1\}$, $L$ is analytic with $L(\rho)=0$ and
$L'(\rho)\neq0$, so $-(L'/L)$ has a simple pole at $\rho$ with residue $1$; multiplying by the
analytic factor $x^{s}/s$ scales the residue to $x^{\rho}/\rho$ and the leading sign gives
$-x^{\rho}/\rho$.
\end{proof}

\section{Projection and reality on the critical line}\label{sec:projection}

The carrier is three-dimensional; the readout projects it to a one-dimensional coordinate. We first
record that coordinate and the change of scale it carries
(Section~\ref{sub:readout}), then the reality of the completed object that the readout returns
(Section~\ref{sub:reality}).

\subsection{The readout coordinate and its scaling}\label{sub:readout}

The carrier is built in its native gauge: the unit is $\pi/3$, and a full angular turn is six such
units, $6\cdot(\pi/3)=2\pi$ (\code{spin\_phasor\_period6}, Theorem~\ref{thm:buckets}). The readout that
sends a three-dimensional harmonic value $t\in\R$ down to a one-dimensional coordinate composes a
Cayley map \cite{Ahlfors1979} with a rescaling to the standard unit interval (\code{HarmonicProjection}):
\[
  \mathrm{toCircleAngle}(t)=2\arctan t\in(-\pi,\pi),\qquad
  \mathrm{toLine}(\theta)=\frac{\theta+\pi}{2\pi}\in(0,1),
\]
and $\mathrm{projection}=\mathrm{toLine}\circ\mathrm{toCircleAngle}$. The second map carries the change
of scale: it sends the full turn $2\pi$ to length $1$, so the six $\pi/3$ buckets of the gauge become
six equal sub-intervals of the readout. Under this rescaling the coordinate is rigid --- the Cayley
map is odd (\code{toCircleAngle\_odd}), the normalization is affine hence midpoint-preserving
(\code{toLine\_midpoint}), and the principal-arc ends $\theta=\pm\pi$ map to the two ends of the
interval (\code{toLine\_neg\_pi}, \code{toLine\_pi}) --- so the centre of the gauge is carried to the
centre of the interval.

\begin{theorem}[Coordinate centre; \code{HarmonicProjection.projection\_midline}]\label{thm:midline}
$\mathrm{projection}(0)$ is the midpoint of the readout interval: the centre of the six gauge buckets,
the angle $\theta=0$, maps to the average of the two endpoint images
(\code{toLine\_midpoint\_of\_endpoints}).
\end{theorem}

The inverse conversion is the same scaling read backwards: a readout coordinate $u$ is the angle
$\theta=2\pi u-\pi$, the gauge reading $6u-3$ in $\pi/3$ units, and the interval's centre is the gauge
centre $\theta=0$. This is a property of the readout coordinate and its scaling. The value the readout
returns on the completed object is real --- it lands on the real axis (Section~\ref{sub:reality}).

\subsection{Reality on the line}\label{sub:reality}

The supporting analytic facts are the conjugation (Schwarz reflection \cite{Ahlfors1979}) and
functional symmetries of $\Lambda$. The $\tfrac12$ in this subsection is the \emph{reflection axis}
$\Ree s=\tfrac12$ --- the fixed line of $s\mapsto 1-s$ on which $\conj s=1-s$
(Section~\ref{sec:notation}) --- \emph{not} the scale-critical exponent of
Section~\ref{sub:scalecrit}: the reality is read off the functional equation, not the area law.

\begin{theorem}[Conjugation symmetry; \code{HelixCollapse.completedRiemannZeta\_conj}]\label{thm:conj}
For all $s\in\C$, $\;\conj{\Lambda(\conj s)}=\Lambda(s)$.
\end{theorem}

\begin{proof}[Proof sketch]
On $\Ree s>1$ this is the real Dirichlet coefficients of $\Lambda$ ($\conj{\Lambda(\conj s)}=\Lambda(s)$
termwise, using $\conj{\pi^{-s/2}}=\pi^{-\conj s/2}$ and $\conj{\Gamma(\conj{s}/2)}=\Gamma(s/2)$); the
identity extends to $\C$ by analytic continuation of $\Lambda_0$ (the pole-removed completion) via the
identity theorem.
\end{proof}

\begin{theorem}[Reality on the line; \code{completedRiemannZeta\_critical\_line\_im\_zero}]\label{thm:real}
For every $t\in\R$, $\;\Imm\,\Lambda(\tfrac12+it)=0$.
\end{theorem}

\begin{proof}[Proof sketch]
On the line $\conj{\tfrac12+it}=1-(\tfrac12+it)$, so Theorem~\ref{thm:conj} together with
$\Lambda(1-s)=\Lambda(s)$ (\code{completedRiemannZeta\_one\_sub}) gives
$\conj{\Lambda(\tfrac12+it)}=\Lambda(\tfrac12+it)$, i.e.\ a real value.
\end{proof}

\begin{theorem}[Faithful projection for $\zeta$; \code{faithful\_projection\_zeta}]\label{thm:faithproj}
For every $\gamma\in\R$,
\[
  \Imm\,\Lambda\!\bigl(\tfrac12+i\gamma\bigr)=0
  \qquad\text{and}\qquad
  F_\eta(\gamma)=0 \iff \zeta\!\bigl(\tfrac12+i\gamma\bigr)=0,
\]
where $F_\eta(\gamma)=(1-2^{1-s})\zeta(s)$ at $s=\tfrac12+i\gamma$ (\code{EtaTrivial.Feta}).
\end{theorem}

\begin{theorem}[Located crossings lie on the line; \code{crossing\_is\_zero\_on\_line}]\label{thm:located}
If the eta-regularized limit of Theorem~\ref{thm:cont} vanishes at ordinate $y$, then
$\zeta(\tfrac12+iy)=0$ and $\Ree(\tfrac12+iy)=\tfrac12$.
\end{theorem}

\section{Faithfulness for every Dirichlet $L$-function and Tate completion}\label{sec:allL}

A single carrier serves all $\chi$; the character weight $\chi(n)$ on each phasor is the only dial.

\begin{theorem}[One carrier, all characters; \code{faithful\_all\_L}]\label{thm:alll}
For every Dirichlet character $\chi$:
\begin{enumerate}[label=\textup{(\roman*)}]
\item for $\Ree s>1$, $\;\sum_{n<N}\chi(n)\,n^{-s}\to L(s,\chi)$; and
\item on the critical line $\Ree s=\tfrac12$, the eta-twisted closed carrier
$\code{etaTwistClosed}(\chi,s)$ vanishes if and only if $L(s,\chi)=0$
(\code{etaTwistClosed\_eq\_zero\_iff\_critical}).
\end{enumerate}
\end{theorem}

\begin{definition}[Completed carrier; \code{Tate.completedCarrier}]\label{def:completed}
For a Dirichlet character $\chi$, set $\;\Lambda_\chi^{\mathrm{c}}(y)=\Lambda(\chi,\tfrac12+iy)$ using
the archimedean $\Gamma$-completion $\Lambda(\chi,s)=(\text{gamma factor})\cdot L(s,\chi)$.
\end{definition}

\begin{theorem}[Completion and functional equation; \code{faithful\_all\_L\_completed}, \code{completedCarrier\_eq\_zero\_iff}, \code{completed\_functional\_equation}]\label{thm:tate}
For every \emph{primitive} $\chi$ modulo $q$:
\begin{enumerate}[label=\textup{(\roman*)}]
\item the completion is zero-free off the zeros of $L$ on the line: $\Lambda_\chi^{\mathrm{c}}(y)=0
\iff L(\tfrac12+iy,\chi)=0$ (the $\Gamma$-factor never vanishes for $\Ree s>0$); and
\item Tate's functional equation \cite{Tate1967} holds:
$\;\Lambda(\chi,1-s)=q^{\,s-1/2}\,W(\chi)\,\Lambda(\chi^{-1},s)$,
where $W(\chi)=\code{rootNumber}(\chi)$.
\end{enumerate}
\end{theorem}

\section{The double-ended carrier and conjugacy}\label{sec:double}

The carrier is chiral: the right strand winds with spin $e^{-iy\log n}$, and its mirror, the left
strand, winds with $e^{+iy\log n}$.

\begin{theorem}[Conjugate chiralities; \code{both\_helices\_conjugate}]\label{thm:conjhelix}
For every $N\in\N$ and $y\in\R$,
\[
  \sum_{n<N}(-1)^{\,n+1}\,n^{-(1/2-iy)}
   \;=\; \conj{\;\sum_{n<N}(-1)^{\,n+1}\,n^{-(1/2+iy)}\;}.
\]
\end{theorem}

\begin{proof}[Proof sketch]
Termwise: $\conj{(-1)^{n+1}}=(-1)^{n+1}$ and $\conj{\,n^{-(1/2+iy)}}=n^{-(1/2-iy)}$ for the positive
real base $n$ (so $\arglocal n\neq\pi$), by $\code{Complex.cpow\_conj}$ and $\conj n=n$.
\end{proof}

\begin{corollary}[Equal modulus, identical crossing heights]\label{cor:equalmod}
The two strands have equal modulus at every $N$, so their limits vanish at the same ordinates $y$.
\end{corollary}

\begin{theorem}[Completed sides agree; \code{both\_sides\_cross\_together}]\label{thm:bothsides}
For every $y\in\R$, $\;\Lambda(\tfrac12-iy)=\Lambda(\tfrac12+iy)$.
\end{theorem}

\begin{proof}[Proof sketch]
$\tfrac12-iy=1-(\tfrac12+iy)$, so this is $\Lambda(1-s)=\Lambda(s)$
(\code{completedRiemannZeta\_one\_sub}) on the line.
\end{proof}

\section{The Frobenius similitude and the determinant identity}\label{sec:det}

The map $n\mapsto pn$ --- and more generally $n\mapsto mn$ --- is the \emph{Frobenius
self-similarity} of the carrier. We define it precisely, as a similitude of the carrier line.

\begin{definition}[Frobenius similitude; \code{frobeniusMultiplier}, \code{frobeniusSimilitude}]\label{def:simil}
For an integer $m$ the \emph{Frobenius multiplier} is
$\mu_\theta(m)=\sqrt m\,\mathrm{wind}_\theta(m)\in\C$, and the \emph{Frobenius similitude}
$\mathrm{Sim}_\theta(m)$ is the $\C$-linear self-map $w\mapsto\mu_\theta(m)\,w$ of the carrier line $\C$.
\end{definition}

\begin{theorem}[Prime multiplication acts by similitude; \code{helixPt\_mul}, \code{helixPt\_eq\_similitude}, \code{norm\_frobeniusMultiplier}, \code{frobeniusSimilitude\_mul}]\label{thm:simil}
For $m,n\neq0$, multiplication $n\mapsto mn$ acts on the carrier point
$\mathrm{helixPt}_\theta(n)=\sqrt n\,\mathrm{wind}_\theta(n)$ by the similitude
\[
  \mathrm{helixPt}_\theta(mn) \;=\; \mathrm{Sim}_\theta(m)\bigl(\mathrm{helixPt}_\theta(n)\bigr)
   \;=\; \underbrace{\sqrt m}_{\text{radial}}\;\underbrace{\mathrm{wind}_\theta(m)}_{\text{rotation}}\;\mathrm{helixPt}_\theta(n).
\]
Its multiplier has modulus $\lVert\mu_\theta(m)\rVert=\sqrt m$ --- the area-law factor of
Theorem~\ref{thm:arealaw}, the winding rotation being unit-modulus --- so the map scales every length
by $\sqrt m$. Moreover the similitudes \emph{compose},
$\mathrm{Sim}_\theta(mn)=\mathrm{Sim}_\theta(m)\circ\mathrm{Sim}_\theta(n)$: the carrier is
\emph{self-similar} under the whole multiplicative monoid generated by the primes (the Frobenius maps).
\end{theorem}

Under the bridge (Theorem~\ref{thm:bridge}), with prime-angle assignment $\theta(p)=-y\log p$ the
rotation eigenphase at ordinate $y$ is the spin $z=\mathrm{spin}(y,n)=e^{-iy\log n}\in\T$
(\code{frobeniusMultiplier\_bridge}). The two chiralities (Theorem~\ref{thm:conjhelix}) carry
$z$ and $\conj z$.

\begin{theorem}[Conjugate-eigenstate determinant; \code{frobenius\_conjugate\_det\_one}]\label{thm:det}
For every $y\in\R$ and $n\in\N$, with $z=\mathrm{spin}(y,n)=e^{-iy\log n}$,
\[
  \det
  \begin{pmatrix} z & 0\\[2pt] 0 & \conj z \end{pmatrix}
  \;=\; z\,\conj z \;=\; |z|^2 \;=\; 1 .
\]
\end{theorem}

\begin{proof}
The $2\times2$ determinant is $z\cdot\conj z-0\cdot0=z\,\conj z$. Since $\conj z=\conj{e^{-iy\log n}}
=e^{\,\conj{-iy\log n}}=e^{+iy\log n}$ and $-iy\log n+\overline{-iy\log n}=0$, we have
$z\,\conj z=e^{-iy\log n}\,e^{+iy\log n}=e^{0}=1$.
\end{proof}

\begin{theorem}[Unitary transverse block; \code{frobeniusBlock\_det\_one}, \code{frobeniusBlock\_unitary}]\label{thm:unitary}
The transverse block $\diag(z,\conj z)$ assembled from the two chiral eigenphases at a conjugate
crossing is \emph{unitary}, $M^{\mathsf H}M=I$, with determinant $z\,\conj z=|z|^2=1$ --- the
orientation- and volume-preserving chiral invariant of the transverse phase dynamics. The radial
scaling $\sqrt m$ of the similitude (Theorem~\ref{thm:simil}) is carried separately and is not part of
this block.
\end{theorem}

\subsection{The self-adjoint eigenstate}\label{sub:eigenstate}

The determinant identity above is stated for the bare eigenphases $z,\conj z\in\T$. We realize them
as the values of a genuine eigenstate of a self-adjoint generator. This is the operator
\emph{design}: it is real because its inputs are real, and it stands independently of any zero.
Reading the nontrivial zeros as the spectrum of a self-adjoint operator is the
\emph{Hilbert--P\'olya} program --- supported by the random-matrix law of the zero statistics
\cite{Montgomery1973} and pursued through the Berry--Keating $H=xp$ quantization
\cite{BerryKeating1999} and Connes's noncommutative-geometry trace formula \cite{Connes1999}; the
operator design here is unconditional and per-height, and makes no claim to be that canonical
operator (Remark~\ref{rem:openlink}).

\begin{definition}[Spectral wave; \code{spectralWave}]\label{def:wave}
For $\gamma\in\R$, $\;\psi_\gamma(t)=e^{i\gamma t}\;(t\in\R)$.
\end{definition}

\begin{theorem}[Unit-modulus eigenstate of $D=-i\,d/dt$; \code{spectralWave\_eigen}, \code{spectralWave\_norm}]\label{thm:wave}
For all $t\in\R$, $\;-i\,\psi_\gamma'(t)=\gamma\,\psi_\gamma(t)\;$ and $\;\|\psi_\gamma(t)\|=1$. Thus
$\psi_\gamma$ is a unit-modulus eigenstate of the generator $D=-i\,d/dt$ with the real eigenvalue
$\gamma$.
\end{theorem}

The reality is structural, not an accident of $\psi_\gamma$. Multiplication by the real height,
$\mathrm{vonNeumannOp}(\gamma)=\gamma\cdot\mathrm{id}$ on the fiber $\C$, is symmetric with
eigenvalue $\gamma$ (\code{vonNeumannOp\_isSymmetric}, \code{vonNeumannOp\_hasEigenvalue}); on the
finitely-supported sequence space the real-diagonal operator $\mathrm{diagOp}(d)$ is symmetric
(\code{diagOp\_symmetric}), with the basis vectors as explicit eigenvectors carrying the real
eigenvalues $d_i$ (\code{diagOp\_eigenvector}), and \emph{every} eigenvalue is real
(\code{diagOp\_spectrum\_real}) by von Neumann reality --- a symmetric operator on a complex
inner-product space has real eigenvalues (\code{symmetric\_eigenvalue\_real}). These are collected in
\code{unconditional\_frobenius\_eigenstate} and \code{frobenius\_spectralWave\_eigenstate}. The
eigenwave is a one-parameter unitary group, $\psi_\gamma(s+t)=\psi_\gamma(s)\,\psi_\gamma(t)$ with
$\psi_\gamma(0)=1$ (\code{spectralWave\_add}, \code{spectralWave\_zero}): the unitary flow generated by
the real spectral parameter $\gamma$ (Stone form).

\begin{proposition}[The chiralities are eigenstate values; \code{spin\_eq\_spectralWave}, \code{conj\_spin\_eq\_spectralWave}]\label{prop:chiral}
For every $y\in\R$ and $n\in\N$, $\;\mathrm{spin}(y,n)=\psi_y(-\log n)\;$ and
$\;\conj{\mathrm{spin}(y,n)}=\psi_y(\log n)$.
\end{proposition}

The useful spectral content is not the mere existence of the eigenwave $\psi_y$ --- every real $y$
gives one --- but its \emph{orthogonality} to the arithmetic data at a vanishing.

\begin{proposition}[The accumulation is an eigenstate--data inner product; \code{fiberCarrierFinite\_eq\_inner}, \code{fiberCarrierFinite\_eq\_zero\_iff\_orthogonal}]\label{prop:ortho}
Split the critical-line phasor $\chi(n)\,n^{-(1/2+iy)}$ into the arithmetic datum
$\chi(n)\,n^{-1/2}$ (the finitely-supported vector $\code{dataVec}(\chi,N)$) and the eigenwave sample
$\conj{\mathrm{spin}(y,n)}=\psi_y(\log n)$ (the vector $\code{waveVec}(y,N)$). With the $\ell^2$ inner
product $\langle f,g\rangle=\sum_n\conj{f(n)}\,g(n)$ on the finitely-supported sequence space, the
finite accumulation is their pairing,
\[
  \sum_{n<N}\chi(n)\,n^{-(1/2+iy)} \;=\; \bigl\langle\,\code{waveVec}(y,N),\ \code{dataVec}(\chi,N)\,\bigr\rangle ,
\]
so it vanishes \emph{iff} the sampled eigenwave is $\ell^2$-orthogonal to the arithmetic data
(\code{fiberCarrierFinite\_eq\_zero\_iff\_orthogonal}). A vanishing point is exactly a real-frequency
orthogonality event, exhibited at each $N$ (and the pairing is the partial phasor sum
$G_{\chi,N}(\tfrac12+iy)$, which accumulates to the carrier value $L(\tfrac12+iy,\chi)$;
\code{inner\_eq\_finiteCarrier}). The Gram form of any finite family of such vectors is
positive-semidefinite (\code{carrier\_gram\_nonneg}), the positive-kernel background for the de Branges
evaluation of Section~\ref{sec:debranges}.
\end{proposition}

\begin{theorem}[Determinant on the eigenstate values; \code{frobenius\_spectralWave\_det\_one}]\label{thm:wavedet}
For every $y\in\R$ and $n\in\N$, $\;\det\diag\!\big(\psi_y(-\log n),\,\psi_y(\log n)\big)=1$.
\end{theorem}

\begin{theorem}[Frobenius unimodularity of the produced eigenstate; \code{frobenius\_eigenstate\_det\_one}]\label{thm:eigdet}
Fix $\gamma\in\R$, $n\in\N$, and $\psi:\R\to\C$, and \emph{assume} the vanishing $\Rightarrow$
eigenstate link $\psi=\psi_\gamma$. Then $\psi$ is a unit-modulus eigenstate of $D=-i\,d/dt$ with
eigenvalue $\gamma$, and $\;\det\diag\!\big(\psi(-\log n),\,\psi(\log n)\big)=1$.
\end{theorem}

\begin{remark}[The open link]\label{rem:openlink}
We exhibit the eigenstate $\psi_\gamma$ for every real $\gamma$ unconditionally
(Theorem~\ref{thm:wave}). What is \emph{not} proved is that a carrier vanishing at height $\gamma$
\emph{produces} this eigenstate; that hypothesis ($\psi=\psi_\gamma$) is the one input of
Theorem~\ref{thm:eigdet}, left open. We do not derive that the vanishings are the spectrum of a
canonical operator.
\end{remark}

\section{The chiral cup and its Hilbert completion}\label{sec:cup}

The pairing underlying the on-line spectral reading is the \emph{chiral cup}, a Hermitian form on the
finitely-supported fibre space $\ell_0=\code{(}\N\to_0\C\code{)}$ assembled from the forward and reverse
(right/left) chiralities of the double-ended carrier (Section~\ref{sec:double}).

\begin{definition}[Chiral defect and cup; \code{CupIdentity.Dop}, \code{Cup}, \code{Jconj}, \code{iotaR}, \code{iotaL}]\label{def:chiralcup}
With the right/left site embeddings $\iota_R,\iota_L$ and the conjugation $J$, the \emph{chiral defect}
is $\mathrm{Dop}\,F=\iota_R F-J(\iota_L F)$ --- the failure of the forward embedding to match the
conjugated reverse --- and the \emph{cup} is the sesquilinear pairing
$\mathrm{Cup}(F,G)=\langle \mathrm{Dop}\,F,\ \mathrm{Dop}\,G\rangle$.
\end{definition}

\begin{theorem}[The cup is a definite Hermitian form; \code{cup\_hermitian}, \code{cup\_positive\_semidefinite}, \code{cup\_self\_im}, \code{cup\_nullspace\_safe}, \code{cup\_null\_iff}]\label{thm:cuppos}
$\mathrm{Cup}$ is Hermitian ($\mathrm{Cup}(G,F)=\conj{\mathrm{Cup}(F,G)}$), positive semidefinite
($0\le\Ree\,\mathrm{Cup}(F,F)$), real on the diagonal ($\Imm\,\mathrm{Cup}(F,F)=0$), and \emph{definite}:
$\mathrm{Cup}(F,F)=0\iff \mathrm{Dop}\,F=0\iff F=0$ (\code{cup\_null\_iff},
\code{dop\_eq\_zero\_iff\_carrier\_zero}). The closure therefore preserves all arithmetic content
(\code{closure\_preserves\_arithmetic\_content}): it never collapses a nonzero fibre, and the readout
$\mathrm{Dop}$ is injective (\code{spectral\_readout\_faithful}).
\end{theorem}

\begin{theorem}[Hilbert completion; \code{hilbert\_completion\_exists}, \code{cup\_self\_eq\_two\_energy}]\label{thm:hilbert}
The cup completes to a genuine Hilbert space: there is a complete inner-product space into which
$\mathrm{Dop}$ embeds with $\langle \mathrm{Dop}\,F,\mathrm{Dop}\,G\rangle=\mathrm{Cup}(F,G)$. On the
diagonal the cup is the $\ell^2$ energy, $\mathrm{Cup}(F,F)=2\sum_n\|F(n)\|^2$.
\end{theorem}

\begin{theorem}[Frobenius transport is a cup similitude; \code{Tw}, \code{transport\_cup\_modulus}, \code{frobenius\_weighted\_cup\_identity}, \code{transport\_unitary\_after\_normalization}]\label{thm:transport}
The weight transport $T_\mu$ scales the cup by $\|\mu\|^2$,
$\langle T_\mu(\mathrm{Dop}\,F),T_\mu(\mathrm{Dop}\,G)\rangle=\|\mu\|^2\,\mathrm{Cup}(F,G)$, and a unit
weight ($\|\mu\|=1$) transports \emph{unitarily} (\code{transport\_unitary\_after\_normalization}). With
$\mu=\mathrm{amp}\cdot\mathrm{sgn}$, $\|\mathrm{sgn}\|=1$, the cup scales by $\mathrm{amp}^2$: this is the
cup-level form of the Frobenius similitude (Theorem~\ref{thm:simil}), whose radial factor $\sqrt m$
squares to the area-law weight $m$.
\end{theorem}

\begin{theorem}[Von Neumann reality and modulus balance; \code{von\_neumann\_reality}, \code{diagOp\_spectrum\_real}, \code{midpoint\_forcing}]\label{thm:vnreality}
Every eigenvalue of a symmetric operator on a complex inner-product space is real
(\code{von\_neumann\_reality}), and the real-diagonal operator $\mathrm{diagOp}\,d$ has real spectrum
(\code{diagOp\_spectrum\_real}). The cup-level reflection pins the readout to the line: for $q>1$, if
$\|q^{s}\|^2=q$ then $\Ree s=\tfrac12$ (\code{midpoint\_forcing}) --- since
$\|q^{s}\|^2=q^{2\Ree s}=q$ forces $2\Ree s=1$, a balance-to-line implication independent of any zero.
\end{theorem}

\begin{theorem}[No chiral dominance, by crossing induction; \code{no\_dominance\_by\_crossing\_induction}, \code{crossing\_balance\_preserved}, \code{crossing\_no\_dominance\_cup}]\label{thm:crossing}
For a crossing helix with a positive Frobenius cup, weighted forward/reverse transport, and
zero-readout exhaustion, the modulus balance $\|\mathrm{fwd}(n)\|=\|\mathrm{rev}(n)\|$ propagates from
each crossing to the next (\code{crossing\_balance\_preserved}); hence no chirality dominates at any
crossing (\code{no\_dominance\_by\_crossing\_induction}) and the forward and reverse cup energies agree
at every site (\code{crossing\_no\_dominance\_cup}). The hypotheses are inhabited and non-vacuous: a
balanced helix satisfies all of them, while an unbalanced one fails exhaustion
(\code{balancedHelix\_exhaust}, \code{unbalancedHelix\_not\_exhaust}).
\end{theorem}

\begin{remark}[The cup carries the von Mangoldt prime field; \code{PrimePowerProbe.cup\_intertwining}, \code{vonMangoldtChannel\_eq\_neg\_logDeriv}]\label{rem:cupvm}
On the prime-power probe fibre the cup energy is exactly the squared norm of the winding-weighted von
Mangoldt terms, $\mathrm{Cup}(\mathrm{probe},\mathrm{probe})=2\sum_{n<N}\|\,\mathrm{channelTerm}(n)\|^2$
(\code{cup\_intertwining}), supported on prime powers (\code{channelTerm\_eq\_zero\_of\_not\_isPrimePow}),
whose superposition is $-L'/L$ for $\Ree s>1$ (\code{vonMangoldtChannel\_eq\_neg\_logDeriv}); in that
region the channel never vanishes (\code{noDominance\_no\_offline\_zero}). This is the explicit formula
of Section~\ref{sec:explicit} read through the cup energy.
\end{remark}

\section{The K\"ahler polarization, the trace formula, and a cohomological re-expression of the zeros}\label{sec:kahler}

The fibre carries a compatible triple --- a symplectic form, a complex structure, and a metric --- under
which the Frobenius maps act as conformal symplectic similitudes, the von Mangoldt prime field appears
as an operator trace, and the vanishing of $L$ is re-expressed as a one-dimensional cohomology.

\subsection{The polarization and its K\"ahler identity}

\begin{definition}[Symplectic form, complex structure, metric; \code{HelixPolarization.OmegaH}, \code{JH}, \code{gH}]\label{def:polarization}
On the fibre $H=\code{(}\N\to_0\C\code{)}$: the real antisymmetric $\R$-bilinear form $\Omega$
(\code{OmegaH\_antisymm}, \code{OmegaH\_add\_right}, \code{OmegaH\_smul\_right}), the almost-complex
structure $J$ with $J^2=-\mathrm{id}$ (\code{JH\_sq}), and the metric $g$.
\end{definition}

\begin{theorem}[K\"ahler compatibility; \code{OmegaH\_JH}, \code{gH\_symm}, \code{OmegaH\_compat}, \code{inner\_eq\_gH\_add\_I\_OmegaH}]\label{thm:kahler}
The triple is compatible: $\Omega(x,Jy)=g(x,y)$, $g$ is symmetric, and $J$ preserves $\Omega$
($\Omega(Jx,Jy)=\Omega(x,y)$). The Hermitian inner product splits into its metric and symplectic parts,
\[
  \langle x,y\rangle \;=\; g(x,y) \;+\; i\,\Omega(x,y),
\]
the K\"ahler identity (\code{inner\_eq\_gH\_add\_I\_OmegaH}).
\end{theorem}

\begin{theorem}[Positivity of the polarization; \code{theoremA}, \code{gH\_Dop\_self}]\label{thm:thmA}
For $F\neq0$, $\;0<\Omega(\mathrm{Dop}\,F,\ J\,\mathrm{Dop}\,F)$: the polarization is strictly positive on
the chiral defect (\code{theoremA}). The metric recovers the cup energy,
$g(\mathrm{Dop}\,F,\mathrm{Dop}\,F)=\Ree\,\mathrm{Cup}(F,F)$ (\code{gH\_Dop\_self}).
\end{theorem}

\subsection{Frobenius as a conformal symplectic similitude}

\begin{theorem}[Frobenius scales the symplectic form by $m$; \code{frobT\_weighted\_adjoint}, \code{theoremB}, \code{frobT\_mul}, \code{siteShift\_weighted\_adjoint}, \code{siteShift\_mul}]\label{thm:thmB}
The Frobenius transport $\mathrm{frobT}(m,\mathrm{sgn})$ (unit sign) scales the inner product by $m$,
$\langle \mathrm{frobT}\,x,\mathrm{frobT}\,y\rangle=m\,\langle x,y\rangle$
(\code{frobT\_weighted\_adjoint}), and on the polarization it scales the symplectic pairing by exactly
$p^{r}$,
\[
  \Omega\bigl(\mathrm{frobT}(p^{r})\,x,\ J\,\mathrm{frobT}(p^{r})\,y\bigr)
   \;=\; p^{r}\,\Omega(x,Jy)\qquad(\text{\code{theoremB}}).
\]
The maps compose multiplicatively, $\mathrm{frobT}(mn,s_1s_2)=\mathrm{frobT}(m,s_1)\circ\mathrm{frobT}(n,s_2)$
(\code{frobT\_mul}). The companion geometric site shift
$\mathrm{siteShift}(m):\delta_n\mapsto\sqrt m\,\delta_{mn}$ is the same $\sqrt m$ similitude on sites,
unitary up to the $m$-factor and multiplicative (\code{siteShift\_weighted\_adjoint},
\code{siteShift\_mul}) --- the operator avatar of $\mathrm{helixPt}(mn)=\sqrt m\,\mathrm{wind}(m)\,\mathrm{helixPt}(n)$.
\end{theorem}

\subsection{The von Mangoldt trace}

\begin{theorem}[The weighted trace is $-L'/L$; \code{theoremC}, \code{traceWeighted\_eq\_superposition}, \code{theoremC\_trace}, \code{specDiag\_siteShift}, \code{localTrace\_mul}]\label{thm:thmC}
The weighted local trace is the prime-power superposition
$\mathrm{traceWeighted}(\sigma,h)=\sum_{n}\chi(n)\Lambda(n)\,\langle\delta_n,\mathrm{specDiag}\,\delta_n\rangle$
(\code{theoremC\_trace}), equal to the von Mangoldt channel (\code{traceWeighted\_eq\_superposition}),
which for $\sigma>1$ is the logarithmic derivative
$-L'(\,\cdot\,)/L(\,\cdot\,)$ at the spectral parameter (\code{theoremC}). The local factor is
completely multiplicative (\code{localTrace\_mul}) and intertwines with the site shift,
$\mathrm{specDiag}\circ\mathrm{siteShift}(m)=\mathrm{localTrace}(m)\cdot(\mathrm{siteShift}(m)\circ\mathrm{specDiag})$
(\code{specDiag\_siteShift}). This is the geometric explicit formula (Section~\ref{sec:explicit})
realized as an operator trace.
\end{theorem}

\subsection{A cohomological re-expression of the vanishing}

\begin{theorem}[$L$ vanishes iff the helix cohomology is nontrivial; \code{theoremD}, \code{theoremD\_H0\_finrank}, \code{theoremD\_H1\_finrank}, \code{theoremD\_diff}]\label{thm:thmD}
The helix cohomology at $\rho$ is the quotient $H^\bullet(\rho)=\C/\langle L(\rho,\chi)\rangle$, with the
eigenclass the class of $1$. Then
\[
  L(\rho,\chi)=0 \;\iff\; \text{eigenclass}(\rho)\neq0 \;\iff\; H^1(\rho)\ \text{nontrivial},
\]
and the cohomology dimension is the exact indicator
$\dim_\C H^0(\rho)=\dim_\C H^1(\rho)=\mathbf 1[\,L(\rho,\chi)=0\,]$ (\code{theoremD\_H0\_finrank},
\code{theoremD\_H1\_finrank}): a simple zero is a one-dimensional class, a non-zero is acyclic.
\end{theorem}

\begin{remark}[The cohomology is the quotient $\C/\langle L\rangle$]\label{rem:cohtaut}
Concretely $H^\bullet(\rho)=\C/\langle L(\rho,\chi)\rangle$, so Theorem~\ref{thm:thmD} is the
field-theoretic identity ``$\C/\langle x\rangle\neq0\iff x=0$'' at $x=L(\rho,\chi)$ --- a cohomological
re-expression of the vanishing. The analytic weight of this section is carried by the polarization
(Theorem~\ref{thm:kahler}), the Frobenius scaling (Theorem~\ref{thm:thmB}), and the trace identity
(Theorem~\ref{thm:thmC}).
\end{remark}

\begin{theorem}[Germ cohomology dimension; \code{germH1\_finrank}, \code{jetRead\_surjective}, \code{ker\_jetRead\_eq\_range\_germCoboundary}]\label{thm:germ}
The $n$-jet germ cohomology has dimension exactly $n$: $\dim_\C H^1_{\mathrm{germ}}(n)=n$
(\code{germH1\_finrank}), via the surjective jet readout (\code{jetRead\_surjective}) and
$\ker(\mathrm{jetRead})=\mathrm{range}(\mathrm{germCoboundary})$
(\code{ker\_jetRead\_eq\_range\_germCoboundary}): the local jet data of order $n$ is an $n$-dimensional
obstruction space.
\end{theorem}

\section{The de Branges evaluation}\label{sec:debranges}

The reality on the critical line (Section~\ref{sub:reality}) and the eigenstate above sit inside the
Hermite--Biehler / de Branges spectral framework \cite{Levin1980,deBranges1968}. We record the framework, its reality and
discreteness consequences, a worked example, and the change of variables that places the on-line
cancellations as real spectral points.

\begin{definition}[Structure function and components; \code{Estar}, \code{Acomp}, \code{Bcomp}, \code{IsHB}]\label{def:hb}
For $E:\C\to\C$, the reflection is $E^*(z)=\conj{E(\conj z)}$, the components are
$A=\tfrac12(E+E^*)$ and $B=\tfrac{i}{2}(E-E^*)$, and $E$ is \emph{Hermite--Biehler} ($\mathrm{IsHB}\,E$)
when $\|E^*(z)\|<\|E(z)\|$ for every $z$ with $\Imm z>0$.
\end{definition}

\begin{proposition}[Gram positivity; \code{gram\_quadratic\_form\_nonneg}]\label{prop:gram}
For vectors $v_i$ in an inner-product space and scalars $c_i$,
$\;\sum_{i,j}\conj{c_i}\,c_j\,\langle v_i,v_j\rangle=\big\langle\sum_i c_i v_i,\sum_j c_j v_j\big\rangle\ge0$
--- the positive-definite reproducing-kernel fact behind $H(E)$.
\end{proposition}

\begin{theorem}[A positive feature map forces Hermite--Biehler; \code{CupIdentity.deBrangesKernel}, \code{deBrangesKernel\_diag}, \code{featureMap\_forces\_HB}, \code{no\_featureMap\_of\_not\_HB}]\label{thm:rkhs}
The de Branges kernel has diagonal
$K_E(w,w)=\bigl(\|E(w)\|^2-\|E^*(w)\|^2\bigr)/(4\pi\,\Imm w)$ (\code{deBrangesKernel\_diag}). If a feature
map $\varphi$ \emph{reproduces} it --- $\langle\varphi(z),\varphi(w)\rangle=K_E(w,z)$ for all $z,w$ ---
then on the open upper half-plane $\|E^*(w)\|\le\|E(w)\|$, i.e.\ the Hermite--Biehler inequality
(\code{featureMap\_forces\_HB}), and the kernel diagonal is nonnegative there
(\code{featureMap\_forces\_diag\_nonneg}). Contrapositively, no positive (reproducing-kernel) feature map
can exist where $\mathrm{HB}$ fails (\code{no\_featureMap\_of\_not\_HB}). Thus $\mathrm{HB}$ is precisely the
condition for the structure kernel to be a genuine inner-product Gram kernel --- the positive-kernel
background (Proposition~\ref{prop:gram}) made into a dichotomy.
\end{theorem}

\begin{theorem}[Positivity forces a real spectrum; \code{hb\_no\_zero\_upper}, \code{Acomp\_zero\_im\_eq\_zero}, \code{Bcomp\_zero\_im\_eq\_zero}]\label{thm:hbreal}
If $\mathrm{IsHB}\,E$ then $E$ has no zeros in the open upper half-plane, and every zero of $A$ and of
$B$ is real.
\end{theorem}

\begin{theorem}[Discreteness; \code{Bcomp\_zeros\_discrete}]\label{thm:hbdisc}
If, in addition, $E$ is entire, then the zeros of $B$ are isolated: a real, discrete spectrum.
\end{theorem}

\begin{proposition}[Paley--Wiener example; \code{paleyWiener\_isHB}, \code{paleyWiener\_Bcomp}]\label{prop:pw}
$E(z)=e^{-iz}$ is Hermite--Biehler, and its $B$-component is $\sin$, whose spectrum is $\{k\pi\}$.
\end{proposition}

\begin{theorem}[The unconditional Paley--Wiener template; \code{paleyWiener\_no\_zero\_upper}, \code{paleyWiener\_balance\_iff\_real}, \code{paleyWiener\_unimodular\_block\_unifies}, \code{paleyWiener\_Bcomp\_zeros\_discrete}]\label{thm:pwtemplate}
For this genuine Hermite--Biehler $E$ all the conditional readings of the paper hold \emph{unconditionally}
and coincide. There are no zeros in the upper half-plane (\code{paleyWiener\_no\_zero\_upper}), the
structure is rigid off the real axis (\code{paleyWiener\_carrierRigid}), and the three readings are
mutually equivalent and equivalent to reality:
\[
  \det\,\mathrm{conjPairBlock}\!\bigl(E(z)/E^*(z)\bigr)=1
  \;\iff\; \|E(z)\|=\|E^*(z)\| \;\iff\; \Imm z=0
\]
(\code{paleyWiener\_unimodular\_block\_unifies}, \code{paleyWiener\_balance\_iff\_real},
\code{paleyWiener\_conjugate\_det\_one\_iff\_reality}). Its $B$-component zeros are real
(\code{paleyWiener\_Bcomp\_zero\_im\_eq\_zero}) and discrete (\code{paleyWiener\_Bcomp\_zeros\_discrete}),
the spectrum $\{k\pi\}$ of $\sin$. This is the unconditional model of exactly what an $\mathrm{HB}$
structure function delivers --- the unitary conjugate-pair block (Theorem~\ref{thm:unitary}) and the
critical-line readout (Corollary~\ref{cor:online}) become a single equivalence; for $\Lambda$ the
identical package is conditional on $\mathrm{IsHB}\,\Lambda$, i.e.\ RH (Remark~\ref{rem:hbscope}).
\end{theorem}

\begin{theorem}[Critical-line change of variables; \code{deBranges\_var\_im}]\label{thm:cov}
For $\rho\in\C$, $\;\Imm\!\big(-i(\rho-\tfrac12)\big)=\tfrac12-\Ree\rho$; so $z=-i(\rho-\tfrac12)$ is
real iff $\Ree\rho=\tfrac12$.
\end{theorem}

\begin{corollary}[On-line cancellations are real spectral points; \code{criticalLine\_deBranges\_real}, \code{zeroHarmonic\_deBranges\_real}]\label{cor:online}
A carrier vanishing at height $\gamma$ sits at $\rho=\tfrac12+i\gamma$, whose de Branges variable
$z=-i(\rho-\tfrac12)=\gamma$ is real. The same-height cancellations we find are real de Branges
spectral points, evaluable in the reality/discreteness framework of
Theorems~\ref{thm:hbreal}--\ref{thm:hbdisc}.
\end{corollary}

\begin{remark}[Scope: $\mathrm{IsHB}\,\Lambda$ is classical open work, not used here]\label{rem:hbscope}
The generic theorems above are conditional on $\mathrm{IsHB}\,E$; we never apply them to a
$\Lambda$-built structure function, and indeed no such structure function is constructed in this work.
That $\mathrm{IsHB}\,\Lambda$ --- no zeros off the critical line (\code{hb\_no\_zero\_upper}) --- is
equivalent to the Riemann Hypothesis is the \emph{classical} de Branges reformulation
\cite{deBranges1968}; discharging it is open work in that program, which this paper neither formalizes,
assumes, nor rests on. What the paper uses is only the \emph{unconditional} on-line content --- the de
Branges variable is real on the critical line (Corollary~\ref{cor:online}), $\mathrm{IsHB}$-free. We make
no claim about off-line zeros, and the paper's open question is projection primacy
(Remark~\ref{rem:admissiblescope}), not $\mathrm{IsHB}\,\Lambda$.
\end{remark}

\section{The harmonic pencil and the focal eigenheight}\label{sec:pencil}

The positive/negative channels of the 3D phasor (Section~\ref{sec:phasor3d}) assemble into a
$2\times2$ determinantal pencil whose rank-drop locates the zeros --- built in the geometric height
$Z=e^{y}>0$ rather than in $\log n$ (the \emph{no-log-trap} design). The readout dictionary is
$t=\log Z$, and the critical readout $\mathrm{criticalReadout}(Z)=\tfrac12+i\log Z$ has real part
$\tfrac12$ for every real $Z$ (\code{criticalReadout\_re}).

\subsection{The scalar channel closure}

\begin{definition}[Signed and unsigned harmonic channels; \code{HarmonicCell.Achan}, \code{Bchan}, \code{Focal.Pchan}, \code{Mchan}]\label{def:channels}
With amplitude $a_n(t)=(\pi/3)\,n^{-1/2}\exp(-it\log n)$, the positive and negative scalar channels
are $P(t)=\sum_{\chi(n)=+1}a_n(t)$ and $M(t)=\sum_{\chi(n)=-1}a_n(t)$. The \emph{signed} harmonic mode
is $B_\chi(Z)=(\pi/3)\,L(\tfrac12+i\log Z,\chi)$ --- the Abel limit of the channel difference $P-M$ ---
and the \emph{unsigned} mode is $A_\chi(Z)=L(\tfrac32+i\log Z,\chi)$, read at the absolute-convergence
abscissa $\sigma=\tfrac32$ where it never vanishes (\code{Achan\_ne\_zero}).
\end{definition}

\begin{theorem}[Channel closure $\iff$ $L$-zero; \code{Focal.channel\_diff\_tendsto}, \code{channel\_closure\_iff\_L\_zero}]\label{thm:channelclosure}
The scalar channel difference is the readout phasor partial sum, with Abel limit
$(\pi/3)\,\Phi_\chi(Z)=(\pi/3)\,L(\tfrac12+i\log Z,\chi)$ (\code{channel\_diff\_tendsto}). Hence the
channel closes --- the difference $P(\log Z)-M(\log Z)$ tends to $0$ --- \emph{iff}
$L(\tfrac12+i\log Z,\chi)=0$ (\code{channel\_closure\_iff\_L\_zero}).
\end{theorem}

\subsection{The determinantal pencil}

\begin{theorem}[Harmonic pencil determinant; \code{HarmonicCell.harmonicPencil\_det}, \code{harmonic\_pencil\_det\_zero\_iff\_L\_zero}]\label{thm:pencildet}
For the diagonal harmonic pencil
\[
  H_\chi=\begin{pmatrix} A_\chi & B_\chi\\[2pt] \mu A_\chi & \lambda B_\chi\end{pmatrix},
  \qquad \det H_\chi=(\lambda-\mu)\,A_\chi B_\chi .
\]
Under admissibility ($A_\chi\neq0$, $\lambda\neq\mu$) the rank-drop $\det H_\chi=0$ is exactly the
scalar $L$-zero closure $L(\tfrac12+i\log Z,\chi)=0$.
\end{theorem}

\begin{theorem}[Gram rank-drop is the same event; \code{gram\_rank\_drop\_iff\_det\_cell\_zero}, \code{gramH\_rank\_drop\_iff\_L\_zero}, \code{helix\_eigenheight\_real}]\label{thm:gram}
For any cell pencil $L^{\mathrm c}$ the Gram matrix $G^{\mathrm c}=(L^{\mathrm c})^{\mathsf H}L^{\mathrm c}$
is Hermitian positive semidefinite with the same kernel as $L^{\mathrm c}$, and in the square case
$\det G^{\mathrm c}=|\det L^{\mathrm c}|^2$; so the Gram rank-drop $\det G^{\mathrm c}_{\chi,h}=0$ is
again exactly the $L$-zero closure (\code{gramH\_rank\_drop\_iff\_L\_zero}). The generalized
eigenheights of such a positive pencil are real Rayleigh quotients
$Z=\langle v,Av\rangle/\langle v,Bv\rangle$ (\code{helix\_eigenheight\_real},
\code{helix\_eigenheight\_rayleigh}), the no-log-trap height diagonal being $Z_n=n$, not $\log n$.
\end{theorem}

\subsection{Admissibility of the source-height readout}

\begin{theorem}[Real source-height $\Rightarrow$ critical line; \code{Admissible.admissible\_iff\_re\_half}, \code{admissible\_real\_height\_implies\_critical\_readout}, \code{nontrivial\_zero\_represented}]\label{thm:admiss}
A complex point is the readout $\mathrm{criticalReadout}(Z)$ of a real source-height $Z$ \emph{iff} it
lies on the critical line $\Ree=\tfrac12$ (\code{admissible\_iff\_re\_half}); the backward direction
realizes it at $Z=e^{\Imm\rho}$. Equivalently an off-line $\rho$ admits \emph{no} real source-height
(\code{offline\_not\_admissible}): the height solving the readout is real iff $\Ree\rho=\tfrac12$
(\code{height\_im}, its imaginary part being $-(\Ree\rho-\tfrac12)$). Consequently every nontrivial
\emph{critical-line} zero is represented by an admissible real source-height Gram eigen-event
$Z=e^{t}>0$ (\code{nontrivial\_zero\_represented}).
\end{theorem}

\begin{remark}[Projection primacy --- the completeness question; \code{ProjectionPrimacy}, \code{projectionPrimacy\_iff\_exhausts}, \code{zero\_source\_admissibility\_of\_projectionPrimacy}]\label{rem:bridge}
The readout reaches \emph{exactly} the critical line for real $Z$ (\code{criticalReadout\_re}), so the
3D carrier's projection cannot land off-axis and cannot host an off-line zero; Theorem~\ref{thm:admiss}
represents every zero \emph{presented on} the line. The \emph{projection-primacy principle}
(\code{ProjectionPrimacy}) asserts that this projection is \emph{complete} --- that its real-height
readout exhausts the nontrivial zeros --- and is proved equivalent to that exhaustion
(\code{projectionPrimacy\_iff\_exhausts}), hence to GRH for $\chi$; conditional on it, \emph{every}
nontrivial zero (with no on-line hypothesis) is a represented admissible eigen-event
(\code{zero\_source\_admissibility\_of\_projectionPrimacy}, \code{RH\_of\_everyZeroAdmissible}). The same
principle is the representability bridge \code{EveryZeroAdmissible}; these two are the \emph{formalized}
faces of the one undecided question (Remark~\ref{rem:admissiblescope}). The de Branges
$\mathrm{IsHB}\,\Lambda$ condition is the same RH content in the classical de Branges form
(Remark~\ref{rem:hbscope}), which this work connects to but neither formalizes nor relies on --- open
work in the de Branges program, not a conditional of this paper. The open question itself is posed,
never assumed: RH fails only if a factor outside the
3D model produces an off-line zero, which the carrier provably cannot. As an unconditional complement, the perturbation (Rouch\'e-type)
stability of the focal zeros is fully proved: at a simple critical-line zero $\gamma$ where the
$L$-channel dominates the bridge error there is a unique nearby focal zero, with the explicit location
bound $|T_d-\gamma|\le 2\varepsilon/a$ (\code{real\_simple\_zero\_location}, \code{real\_simple\_zero\_exists},
\code{real\_simple\_zero\_unique}).
\end{remark}

\section{Admissibility: the carrier has no native off-axis support}\label{sec:admissible}

The constructions above combine into a single intrinsic, on-axis statement about the carrier's
\emph{state space} --- sharper than ``a phasor plot that finds the zeros'': it reframes the critical
line as the carrier's admissible state space, the object whose completeness is projection primacy
(Remark~\ref{rem:admissiblescope}). A point of the vertical line $s=\sigma+iy$ rides the carrier as a stable,
scale-normalized fiber only when its amplitude $n^{-\sigma}$ reciprocates the emergent area-law radius
$r_n\sim\sqrt n$ to a finite positive scale. We call such a $\sigma$ an \emph{admissible
(scale-balanced) state}.

\begin{definition}[Scale-balanced admissible state; \code{ScaleBalanced}]\label{def:admissible}
For a carrier $(p,r)$ and abscissa $\sigma\in\R$, $\sigma$ is \emph{scale-balanced} when the
scale-balance product has a positive finite limit: $\exists L>0,\ n^{-\sigma}\,r_n\to L$, where
$r_n=\mathrm{carrierRadius}(p,r,n)$ is the emergent area-law radius (Definition~\ref{def:wound},
Theorem~\ref{thm:arealaw}).
\end{definition}

\begin{theorem}[The admissible exponent is forced, not chosen; \code{scaleBalanced\_iff}, \code{scaleBalanced\_setOf\_eq}]\label{thm:admissible}
For every carrier $(p,r)$ with $r>0$, $\sigma$ is scale-balanced \emph{iff} $\sigma=\tfrac12$;
equivalently $\{\sigma\in\R:\sigma\text{ scale-balanced}\}=\{\tfrac12\}$. The critical abscissa is the
\emph{unique} exponent compatible with the arclength area law (Theorem~\ref{thm:scalecrit}), and the
conclusion is gauge-independent: the gauge $(p,r)$ fixes only the value of the limit, never which
$\sigma$ is admissible.
\end{theorem}

\begin{corollary}[No native off-axis support; \code{not\_scaleBalanced\_of\_ne}]\label{cor:offaxis}
For $\sigma\neq\tfrac12$ the fiber is \emph{not} a scale-balanced state: its amplitude $n^{-\sigma}$ is
not scale-normalized to the carrier (it decays faster than the radius grows, or outruns it). The
carrier supports no admissible off-axis state.
\end{corollary}

\begin{remark}[Answering the ``designed around $\tfrac12$'' objection]\label{rem:notchosen}
The exclusion of off-axis states is not built in: $\tfrac12$ is \emph{derived} from the arclength area
law (Theorems~\ref{thm:arealaw},~\ref{thm:scalecrit}) as the unique exponent at which $n^{-\sigma}$
reciprocates $r_n\sim\sqrt n$. It is the $\tfrac12$ of $\sqrt n=n^{1/2}$, not a posited abscissa.
\end{remark}

The functional equation supplies the global counterpart: the reflection $s\mapsto1-s$, which pairs the
two chiral fibers, fixes exactly the critical line.

\begin{theorem}[The reflection fixes exactly the critical line; \code{reflection\_fixes\_iff}]\label{thm:reflfix}
For $s\in\C$, $\;\Ree(1-s)=\Ree s\iff\Ree s=\tfrac12$. Off-axis, $s\mapsto1-s$ sends an abscissa
imbalance $\sigma$ to the opposite imbalance $1-\sigma$ --- a dual \emph{pair}, not a self-conjugate
state --- so the double-ended carrier closes only on the fixed locus $\Ree s=\tfrac12$.
\end{theorem}

\begin{theorem}[The hierarchy, assembled; \code{no\_native\_offaxis\_support}]\label{thm:hierarchy}
For every carrier $(p,r)$ with $r>0$, height $y$, and site $n$, the following hold simultaneously:
\begin{enumerate}[label=\textup{(\arabic*)}]
\item the area law $r_n^{\,2}/n\to r(\pi/3)/\pi$ (so $r_n\sim\sqrt n$), from the arclength geometry;
\item scale balance holds iff $\sigma=\tfrac12$ (forced, not chosen);
\item every $\sigma\neq\tfrac12$ is inadmissible;
\item at the balanced axis the chiral fibers $z,\conj z$ form a determinant-one \emph{unitary}
transverse block (Theorem~\ref{thm:unitary}); and
\item $s\mapsto1-s$ fixes exactly $\Ree s=\tfrac12$, on which the de Branges variable is real
(Corollary~\ref{cor:online}).
\end{enumerate}
The chain
\[
  r_n\sim\sqrt n \;\Longrightarrow\; \sigma=\tfrac12 \;\Longrightarrow\;
  \text{unitary chiral block} \;\Longrightarrow\; \det=1 \;\Longrightarrow\;
  s\leftrightarrow1-s \text{ closes on } \Ree s=\tfrac12
\]
is the carrier's intrinsic logic: it reads as a geometry, not a graphing of $L(\tfrac12+it)$.
\end{theorem}

\begin{remark}[Projection primacy: the pivotal hypothesis]\label{rem:admissiblescope}
Theorem~\ref{thm:hierarchy} establishes that the carrier's admissible state space --- the image of its
real-height projection $Z\mapsto\tfrac12+i\tau_\chi(\log Z)$ --- is the critical line: off-axis points are
unsupported / inadmissible / not scale-normalized, so the 3D representation \emph{cannot} host an off-line
zero (\code{criticalReadout\_re}). Whether the projection is also \emph{complete} --- whether it exhausts
\emph{every} nontrivial zero --- is the \emph{projection-primacy principle} (\code{ProjectionPrimacy}),
proved equivalent to the real-height exhaustion of the zeros (\code{projectionPrimacy\_iff\_exhausts}) and
hence to GRH for $\chi$; conditional on it the exhaustion is unconditional
(\code{zero\_source\_admissibility\_of\_projectionPrimacy}). This is the open consideration, and it is a
genuine dichotomy: \emph{either} the projection is primary --- the more accurate 3D model defines the
1D/2D readout completely, and there are no off-line zeros --- \emph{or} some factor outside the 3D
representation produces an off-line zero, one the carrier provably cannot. The Riemann Hypothesis is thus
not something weaker that this work falls short of, but exactly the completeness of the projection; the
work poses the dichotomy and proves the off-axis branch impossible \emph{within} the model. ``Frobenius''
here means precisely the prime-multiplication similitude of Definition~\ref{def:simil}.
\end{remark}

\section{Formalization}\label{sec:formal}

Every definition and theorem above is formalized in Lean~4 \cite{Lean4} against Mathlib
\cite{Mathlib2020}, in the package
\code{RequestProject} (files \code{ClosedForm.lean}, \code{HelixLogFreeFTA.lean},
\code{HelixCollapseReality.lean}, \code{LFunctionPhasor.lean}, \code{GeometricProjectionHolds.lean},
\code{Faithfulness.lean}, \code{AreaLaw.lean}, \code{UnconditionalFrobenius.lean},
\code{DeBranges.lean}, \code{FrobeniusSimilitude.lean}, \code{PerHeightConvergence.lean},
\code{HeightGrowthActive.lean}, \code{CircleMonodromy.lean}, \code{CriticalLineBridge.lean},
\code{GeometricPhasorClosure.lean}, \code{Phasor3D.lean}, \code{FocalEigenheight.lean},
\code{HarmonicPencilCell.lean}, \code{AdmissibleEigenheight.lean}, \code{FullFiber.lean},
\code{PrimePowerProbe.lean}, \code{CupIdentity.lean}, \code{ChiralCup.lean},
\code{CrossingInduction.lean}, \code{CrossingEncodingAudit.lean}, \code{HelixPolarization.lean},
\code{ConditionalToUnconditional.lean}). Each named
theorem compiles with no \code{sorry}; the kernel axiom footprint reported by \code{\#print axioms}
(equivalently \code{lean\_verify}) is
\[
  \{\ \mathrm{propext},\ \ \mathrm{Classical.choice},\ \ \mathrm{Quot.sound}\ \},
\]
the three standard Mathlib axioms, with no construction-specific axiom. The correspondence between
the statements here and the Lean declarations is given by the identifiers in parentheses throughout.

The capstone identities --- the strip extension (Theorems~\ref{thm:dstrip},~\ref{thm:etastrip},~\ref{thm:cont}),
the conjugacy of the two chiralities (Theorem~\ref{thm:conjhelix}), and the determinant identity
(Theorem~\ref{thm:det}) --- are collected in \code{Faithfulness.lean}; the geometry
(Section~\ref{sec:carrier}) and the log-free winding (Section~\ref{sec:winding}) live in
\code{ClosedForm.lean} and \code{HelixLogFreeFTA.lean}, with the emergent area-law radius
(Theorem~\ref{thm:arealaw}, Definition~\ref{def:wound}) and the scale-critical exponent
(Theorem~\ref{thm:scalecrit}) in \code{AreaLaw.lean}. The self-adjoint eigenstate design
(Section~\ref{sub:eigenstate}) lives in \code{UnconditionalFrobenius.lean} and the Hermite--Biehler /
de Branges framework (Section~\ref{sec:debranges}) in \code{DeBranges.lean}, with their connection to
the carrier --- the eigenstate identification (\code{spin\_eq\_spectralWave},
\code{frobenius\_eigenstate\_det\_one}) and the on-line de Branges evaluation
(\code{criticalLine\_deBranges\_real}, \code{zeroHarmonic\_deBranges\_real}) --- collected in
\code{Faithfulness.lean}. The compact local spectral mechanism of Section~\ref{sec:det} --- the
Frobenius similitude (Definition~\ref{def:simil}, Theorem~\ref{thm:simil}), the unitary transverse
block (Theorem~\ref{thm:unitary}), the eigenstate--data orthogonality
(Proposition~\ref{prop:ortho}), and the admissibility / no-off-axis-support results of
Section~\ref{sec:admissible} (Theorems~\ref{thm:admissible},~\ref{thm:reflfix},~\ref{thm:hierarchy})
--- lives in \code{FrobeniusSimilitude.lean}, assembled in
\code{frobenius\_local\_spectral\_mechanism} and \code{no\_native\_offaxis\_support}.
The remaining layers map as follows: the per-height convergence on the line
(Section~\ref{sub:perheight}) is \code{PerHeightConvergence.lean}; the constant height law and active
fraction (Remark~\ref{rem:height}) and the exact $\pi/3$ monodromy (Theorem~\ref{thm:monodromy}) are
\code{HeightGrowthActive.lean} and \code{CircleMonodromy.lean}; the geometric phasor-closure dictionary
and lane split (Section~\ref{sec:strip}) are \code{GeometricPhasorClosure.lean} and
\code{CriticalLineBridge.lean}; the three-channel 3D phasor (Section~\ref{sec:phasor3d}) is
\code{Phasor3D.lean}; and the harmonic pencil, focal eigenheight, and source-height admissibility
(Section~\ref{sec:pencil}) are \code{HarmonicPencilCell.lean}, \code{FocalEigenheight.lean}, and
\code{AdmissibleEigenheight.lean}. The chiral cup and its Hilbert completion (Section~\ref{sec:cup})
live in \code{CupIdentity.lean}, \code{ChiralCup.lean}, \code{CrossingInduction.lean} (with the
non-vacuity audit \code{CrossingEncodingAudit.lean}) and \code{PrimePowerProbe.lean}; the K\"ahler
polarization, the trace formula, and the cohomological re-expression of the zeros (Section~\ref{sec:kahler}) are
\code{HelixPolarization.lean}; the convergent $\eta$-fibre (Section~\ref{sub:fullfiber}) is
\code{FullFiber.lean}; and the unconditional Paley--Wiener template (Theorem~\ref{thm:pwtemplate}) is
\code{ConditionalToUnconditional.lean}.

\begin{remark}[Scope]\label{rem:scope}
This paper presents a geometric state space: the double-ended Frobenius helix carrier and its phasor
fiber, whose admissible scale-balanced states exist only on the critical line
(Section~\ref{sec:admissible}) and on which the vanishing points are \emph{exactly} the on-line zeros
of $\zeta$ and every Dirichlet $L$-function, by the formalized equivalence (Theorem~\ref{thm:alll},
Section~\ref{sub:fullfiber}). Its central
contribution is intrinsic: the carrier has \emph{no native off-axis support}
(Theorem~\ref{thm:hierarchy}), and RH is reframed as the \emph{projection-primacy} principle --- the
completeness of the carrier's real-height projection (\code{projectionPrimacy\_iff\_exhausts}), equivalent
to GRH, with the off-axis branch proved impossible inside the model.
Its mathematical content is exactly the
statements proved above --- the carrier, the area-law derivation of the critical exponent
(Theorem~\ref{thm:scalecrit}), the accumulation to $L(s,\chi)$ on $\Ree s>0$, the
explicit-formula reading of the zeros as poles, the reality of the completed object on the critical
line, the conjugacy of the two chiralities, the Frobenius similitude (Theorem~\ref{thm:simil}), the
unitary transverse block (Theorem~\ref{thm:unitary}), the eigenstate--data orthogonality
(Proposition~\ref{prop:ortho}), the self-adjoint eigenstate realization of the chiral eigenphases
(Section~\ref{sub:eigenstate}), the de Branges evaluation placing the on-line cancellations as real
spectral points (Section~\ref{sec:debranges}), and the admissibility restriction of the carrier's
state space to the critical line (Section~\ref{sec:admissible}). The pivotal hypothesis is
\emph{projection primacy} --- the completeness of the 3D-to-1D projection
(Remark~\ref{rem:admissiblescope}) --- formalized in two guises: that the real-height readout exhausts
the zeros (\code{projectionPrimacy\_iff\_exhausts}), and that a vanishing \emph{produces} the eigenstate
(Remark~\ref{rem:openlink}). The classical de Branges $\mathrm{IsHB}\,\Lambda$ condition
(Remark~\ref{rem:hbscope}) is the same RH content in the de Branges form --- connected here but neither
formalized nor relied on (open work, not a conditional of this paper). This work does not decide projection primacy: it poses the dichotomy ---
the projection is primary, so there are no off-line zeros, \emph{or} an unmodeled factor produces one ---
and proves that the off-axis branch cannot arise within the 3D carrier. The Riemann Hypothesis is exactly
the resolution of this dichotomy, which this work frames but neither assumes nor proves.
\end{remark}

\begin{thebibliography}{9}

\bibitem{Riemann1859} B.~Riemann,
\emph{Ueber die Anzahl der Primzahlen unter einer gegebenen Gr\"osse},
Monatsberichte der Berliner Akademie, 1859.

\bibitem{Edwards1974} H.~M.~Edwards,
\emph{Riemann's Zeta Function}, Academic Press, 1974.

\bibitem{Titchmarsh1986} E.~C.~Titchmarsh,
\emph{The Theory of the Riemann Zeta-function}, 2nd ed.\ (rev.\ Heath-Brown),
Oxford University Press, 1986.

\bibitem{Davenport2000} H.~Davenport,
\emph{Multiplicative Number Theory}, 3rd ed., Springer GTM~74, 2000.

\bibitem{IwaniecKowalski2004} H.~Iwaniec and E.~Kowalski,
\emph{Analytic Number Theory}, AMS Colloquium Publications~53, 2004.

\bibitem{Apostol1976} T.~M.~Apostol,
\emph{Introduction to Analytic Number Theory}, Springer Undergraduate Texts in Mathematics, 1976.

\bibitem{MontgomeryVaughan2007} H.~L.~Montgomery and R.~C.~Vaughan,
\emph{Multiplicative Number Theory~I: Classical Theory},
Cambridge Studies in Advanced Mathematics~97, 2007.

\bibitem{IrelandRosen1990} K.~Ireland and M.~Rosen,
\emph{A Classical Introduction to Modern Number Theory}, 2nd ed., Springer GTM~84, 1990.

\bibitem{Ahlfors1979} L.~V.~Ahlfors,
\emph{Complex Analysis}, 3rd ed., McGraw-Hill, 1979.

\bibitem{Levin1980} B.~Ya.~Levin,
\emph{Distribution of Zeros of Entire Functions}, rev.\ ed., Translations of Mathematical
Monographs~5, American Mathematical Society, 1980.

\bibitem{deBranges1968} L.~de~Branges,
\emph{Hilbert Spaces of Entire Functions}, Prentice-Hall, Englewood Cliffs, N.J., 1968.

\bibitem{Tate1967} J.~T.~Tate,
\emph{Fourier analysis in number fields and Hecke's zeta-functions}, in
J.~W.~S.~Cassels and A.~Fr\"ohlich (eds.), \emph{Algebraic Number Theory},
Academic Press, 1967, pp.~305--347.

\bibitem{Montgomery1973} H.~L.~Montgomery,
\emph{The pair correlation of zeros of the zeta function}, in \emph{Analytic Number Theory}
(Proc.\ Sympos.\ Pure Math.~XXIV, St.~Louis, 1972), American Mathematical Society, 1973,
pp.~181--193.

\bibitem{BerryKeating1999} M.~V.~Berry and J.~P.~Keating,
\emph{The Riemann zeros and eigenvalue asymptotics}, SIAM Review~\textbf{41} (1999), 236--266.

\bibitem{Connes1999} A.~Connes,
\emph{Trace formula in noncommutative geometry and the zeros of the Riemann zeta function},
Selecta Math.\ (N.S.)~\textbf{5} (1999), 29--106.

\bibitem{Mathlib2020} The mathlib Community,
\emph{The Lean Mathematical Library},
Proceedings of CPP 2020, ACM, 2020.

\bibitem{Lean4} L.~de~Moura and S.~Ullrich,
\emph{The Lean 4 Theorem Prover and Programming Language},
Proceedings of CADE-28, Springer LNCS~12699, 2021.

\end{thebibliography}

\end{document}
